\documentclass[11pt,a4paper]{article}

% ==================== TEMEL PAKETLER (pdflatex) ====================
\usepackage[T1]{fontenc}
\usepackage[utf8]{inputenc}
\usepackage{lmodern}
\usepackage[a4paper,margin=2cm,top=2.4cm,bottom=2.4cm]{geometry}
\usepackage{xcolor}
\usepackage{colortbl}
\usepackage{array}
\usepackage{booktabs}
\usepackage{longtable}
\usepackage{tabularx}
\usepackage{listings}
\usepackage{fancyhdr}
\usepackage{amsmath}
\usepackage{hyperref}

% ==================== RENK PALETİ ====================
\definecolor{anaMavi}{HTML}{1B3A5B}
\definecolor{vurguMavi}{HTML}{2E6DA4}
\definecolor{acikMavi}{HTML}{EAF2F9}
\definecolor{turuncu}{HTML}{D9701E}
\definecolor{yesil}{HTML}{2E7D32}
\definecolor{acikYesil}{HTML}{E8F5E9}
\definecolor{kirmizi}{HTML}{C0392B}
\definecolor{mor}{HTML}{6C3483}
\definecolor{acikMor}{HTML}{F3E9F7}
\definecolor{acikGri}{HTML}{F4F5F7}
\definecolor{ortaGri}{HTML}{6B7280}
\definecolor{koyuGri}{HTML}{2D2D2D}
\definecolor{kodArka}{HTML}{FBFBFA}

% ==================== HYPERREF ====================
\hypersetup{
  colorlinks=true,
  linkcolor={[HTML]{2E6DA4}},
  urlcolor={[HTML]{2E6DA4}},
  pdftitle={STM32 Gömülü Sistemler ve ARM Assembly Rehberi},
  pdfauthor={Gömülü Rehber},
  bookmarksnumbered=true
}

% ==================== BAŞLIK STİLLERİ ====================
\setcounter{secnumdepth}{0}
\newcounter{secc}
\newcounter{subsecc}[secc]
\newcounter{subsubsecc}[subsecc]
\newcommand{\gsec}[1]{\par\phantomsection\stepcounter{secc}%
  \setcounter{subsecc}{0}%
  \vspace{18pt}{\color{anaMavi}\Large\bfseries\thesecc.\hspace{0.5em}#1}\par
  \vspace{2pt}{\color{turuncu}\rule{\linewidth}{1.6pt}}\par\vspace{8pt}%
  \addcontentsline{toc}{section}{\protect\numberline{\thesecc}#1}\markboth{#1}{#1}}
\newcommand{\gsub}[1]{\par\phantomsection\stepcounter{subsecc}%
  \setcounter{subsubsecc}{0}%
  \vspace{11pt}{\color{vurguMavi}\large\bfseries\thesecc.\thesubsecc\hspace{0.5em}#1}\par\vspace{4pt}%
  \addcontentsline{toc}{subsection}{\protect\numberline{\thesecc.\thesubsecc}#1}}
\newcommand{\gsubsub}[1]{\par\phantomsection\stepcounter{subsubsecc}%
  \vspace{7pt}{\color{koyuGri}\normalsize\bfseries #1}\par\vspace{2pt}}

% Kısım (Part) başlığı
\newcommand{\gpart}[2]{\clearpage
  \begin{center}\vspace*{1.5cm}
  {\color{turuncu}\rule{0.7\linewidth}{1.5pt}}\\[10pt]
  {\color{ortaGri}\Large KISIM #1}\\[8pt]
  {\color{anaMavi}\Huge\bfseries #2}\\[10pt]
  {\color{turuncu}\rule{0.7\linewidth}{1.5pt}}
  \end{center}\vspace{1cm}}

% ==================== BAŞLIK / ALTBİLGİ ====================
\setlength{\headheight}{14pt}
\pagestyle{fancy}
\fancyhf{}
\renewcommand{\headrulewidth}{0.4pt}
\renewcommand{\footrulewidth}{0.4pt}
\fancyhead[L]{\small\color{ortaGri}STM32 Gömülü \& ARM Assembly}
\fancyhead[R]{\small\color{ortaGri}\leftmark}
\fancyfoot[C]{\small\color{ortaGri}\thepage}

% ==================== KOD LİSTELEME (C - gömülü) ====================
\lstdefinestyle{cstyle}{
  language=C,
  basicstyle=\ttfamily\small\color{koyuGri},
  keywordstyle=\color{vurguMavi}\bfseries,
  commentstyle=\color{yesil}\itshape,
  stringstyle=\color{turuncu},
  numberstyle=\tiny\color{ortaGri},
  numbers=left,
  numbersep=8pt,
  backgroundcolor=\color{kodArka},
  frame=single,
  framesep=6pt,
  rulecolor=\color{ortaGri!40},
  breaklines=true,
  showstringspaces=false,
  tabsize=4,
  columns=fullflexible,
  keepspaces=true,
  extendedchars=true,
  literate={ı}{{\i}}1{İ}{{\.I}}1{ş}{{\c s}}1{Ş}{{\c S}}1{ğ}{{\u g}}1{Ğ}{{\u G}}1{ç}{{\c c}}1{Ç}{{\c C}}1{ö}{{\"o}}1{Ö}{{\"O}}1{ü}{{\"u}}1{Ü}{{\"U}}1,
  morekeywords={uint8_t,uint16_t,uint32_t,int32_t,volatile,GPIO_TypeDef,
    RCC_TypeDef,TIM_TypeDef,USART_TypeDef,__IO,static,inline,void,
    HAL_StatusTypeDef,uint64_t,size_t,bool,SysTick,NVIC_Type}
}
\lstset{style=cstyle}

% Assembly listeleme stili (düz - dile bağlı hata riskini önler)
\lstdefinestyle{asmstyle}{
  basicstyle=\ttfamily\small\color{koyuGri},
  commentstyle=\color{yesil}\itshape,
  backgroundcolor=\color{kodArka},
  frame=single,framesep=6pt,rulecolor=\color{ortaGri!40},
  numberstyle=\tiny\color{ortaGri},numbers=left,numbersep=8pt,
  breaklines=true,showstringspaces=false,tabsize=8,columns=fullflexible,
  keepspaces=true,extendedchars=true,
  morecomment=[l]{;},
  literate={ı}{{\i}}1{İ}{{\.I}}1{ş}{{\c s}}1{Ş}{{\c S}}1{ğ}{{\u g}}1{Ğ}{{\u G}}1{ç}{{\c c}}1{Ç}{{\c C}}1{ö}{{\"o}}1{Ö}{{\"O}}1{ü}{{\"u}}1{Ü}{{\"U}}1,
}

% Kabuk/terminal stili
\lstdefinestyle{shstyle}{
  basicstyle=\ttfamily\small\color{koyuGri},
  backgroundcolor=\color{koyuGri!6},
  frame=single,framesep=6pt,rulecolor=\color{ortaGri!40},
  breaklines=true,showstringspaces=false,tabsize=4,columns=fullflexible,
  keepspaces=true,
  literate={ı}{{\i}}1{İ}{{\.I}}1{ş}{{\c s}}1{Ş}{{\c S}}1{ğ}{{\u g}}1{Ğ}{{\u G}}1{ç}{{\c c}}1{Ç}{{\c C}}1{ö}{{\"o}}1{Ö}{{\"O}}1{ü}{{\"u}}1{Ü}{{\"U}}1,
}

% ==================== ÖZEL KUTULAR ====================
\newcommand{\callout}[4]{\par\smallskip\noindent\begingroup
  \setlength{\fboxsep}{8pt}\setlength{\fboxrule}{0.8pt}%
  \fcolorbox{#1}{#2}{\parbox{\dimexpr\linewidth-2\fboxsep-2\fboxrule\relax}{%
    {\bfseries\color{#1}#3}\par\smallskip #4}}\endgroup\par\smallskip}
\newcommand{\bilgi}[1]{\callout{vurguMavi}{acikMavi}{Bilgi}{#1}}
\newcommand{\ipucu}[1]{\callout{yesil}{acikYesil}{İpucu}{#1}}
\newcommand{\uyari}[1]{\callout{kirmizi}{kirmizi!7}{Dikkat}{#1}}
\newcommand{\ornek}[2][Örnek]{\callout{ortaGri!90!black}{acikGri}{#1}{#2}}
\newcommand{\notk}[2][Not]{\callout{mor}{acikMor}{#1}{#2}}

% Yardımcı komutlar
\newcommand{\code}[1]{\texttt{\color{koyuGri}#1}}
\renewcommand{\arraystretch}{1.35}

% ==================== BAŞLANGIÇ ====================
\begin{document}

% ---------- KAPAK ----------
\begin{titlepage}
\centering
\vspace*{1.4cm}
{\color{turuncu}\rule{\textwidth}{2pt}}\\[8pt]
{\color{anaMavi}\Huge\bfseries STM32 Gömülü Sistemler}\\[10pt]
{\color{vurguMavi}\LARGE ve ARM Assembly Rehberi}\\[8pt]
{\color{ortaGri}\large Register Seviyesinden Assembly'ye Uygulamalı Başvuru}\\[6pt]
{\color{turuncu}\rule{\textwidth}{2pt}}\\[1.4cm]

\begin{center}
\fcolorbox{vurguMavi}{acikMavi}{\parbox{0.86\textwidth}{\centering\large
Bu belge, STM32 (ARM Cortex-M) mikrodenetleyicilerini \textbf{register
seviyesinden} başlayarak, çevre birimleri, kesmeler ve DMA ile
\textbf{çalışan kod örnekleriyle} anlatır; ardından donanımın altındaki
\textbf{ARM Cortex-M assembly} dilini komut komut öğretir.\vspace{4pt}}}
\end{center}

\vspace{0.7cm}
\begin{center}
\fcolorbox{ortaGri!50}{acikGri}{\parbox{0.86\textwidth}{\small
\textbf{Kısım 1 --- STM32 \& Cortex-M Temelleri:} mimari \textbullet\ bellek
haritası \textbullet\ saat (clock) sistemi \textbullet\ CMSIS \textbullet\ araç zinciri.\vspace{3pt}

\textbf{Kısım 2 --- GPIO ve Register Programlama:} register erişimi \textbullet\
bit işlemleri \textbullet\ LED yakma \textbullet\ buton okuma \textbullet\ bare-metal vs HAL.\vspace{3pt}

\textbf{Kısım 3 --- Çevre Birimleri:} zamanlayıcı (timer) \& PWM \textbullet\
UART \textbullet\ ADC \textbullet\ SPI \& I2C.\vspace{3pt}

\textbf{Kısım 4 --- Kesmeler:} NVIC \textbullet\ EXTI \textbullet\ SysTick
\textbullet\ kesme öncelikleri \textbullet\ DMA.\vspace{3pt}

\textbf{Kısım 5 --- ARM Assembly Temelleri:} Cortex-M programlama modeli
\textbullet\ register'lar \textbullet\ Thumb-2 \textbullet\ temel komutlar.\vspace{3pt}

\textbf{Kısım 6 --- ARM Assembly İleri:} akış kontrolü \textbullet\ fonksiyonlar
\& yığın (AAPCS) \textbullet\ diziler \textbullet\ C ile karıştırma.\vspace{3pt}}}
\end{center}

\vfill
{\color{ortaGri}\large Hazırlanma Tarihi: 13 Temmuz 2026}\\[4pt]
{\color{ortaGri} STM32 \textbullet\ ARM Cortex-M \textbullet\ Kodlar İngilizce, açıklamalar Türkçe}
\vspace*{0.6cm}
\end{titlepage}

% ---------- İÇİNDEKİLER ----------
\tableofcontents
\thispagestyle{empty}
\clearpage

% #####################################################################
\gpart{1}{STM32 ve Cortex-M Temelleri}
% #####################################################################

% =====================================================================
\gsec{Giriş: Mikrodenetleyiciler ve STM32}
% =====================================================================

Mikrodenetleyici (MCU), bir işlemci çekirdeği, bellek (Flash + RAM) ve çevre
birimlerini (GPIO, timer, UART\dots) tek bir çip üzerinde toplayan küçük bir
bilgisayardır. STM32, STMicroelectronics'in ARM \emph{Cortex-M} çekirdeği
kullanan popüler MCU ailesidir; endüstride, robotikte ve hobide yaygındır.

\gsub{Mikroişlemci vs Mikrodenetleyici}

\begin{center}
\begin{tabularx}{\textwidth}{@{}>{\bfseries\color{anaMavi}}l X X@{}}
\toprule
\rowcolor{acikMavi}\color{anaMavi}\textbf{} & \color{anaMavi}\textbf{Mikroişlemci (PC/telefon)} & \color{anaMavi}\textbf{Mikrodenetleyici (STM32)}\\
\midrule
Bellek & Harici GB'larca RAM & Dahili KB'larca RAM/Flash\\
\rowcolor{acikGri} İşletim sistemi & Linux/Windows & Genelde yok (bare-metal) veya RTOS\\
Çevre birimleri & Harici çipler & Çip içinde tümleşik\\
\rowcolor{acikGri} Güç & Watt'lar & Miliwatt'lar\\
Gerçek zaman & Zor (OS gecikmesi) & Deterministik, kolay\\
\bottomrule
\end{tabularx}
\end{center}

\gsub{STM32 Ailesi ve Cortex-M Çekirdekleri}

STM32 ailesi, ihtiyaca göre farklı Cortex-M çekirdekleri kullanır. Hepsi aynı
komut setini (Thumb-2) paylaşır, bu yüzden bilgiler taşınabilirdir.

\begin{center}
\begin{tabularx}{\textwidth}{@{}>{\bfseries\color{anaMavi}}l l X@{}}
\toprule
\rowcolor{acikMavi}\color{anaMavi}\textbf{Seri} & \color{anaMavi}\textbf{Çekirdek} & \color{anaMavi}\textbf{Kullanım}\\
\midrule
STM32F0/L0 & Cortex-M0/M0+ & Düşük maliyet, düşük güç\\
\rowcolor{acikGri} STM32F1/F3 & Cortex-M3/M4 & Genel amaçlı (ör. ünlü ``Blue Pill'')\\
STM32F4/F7 & Cortex-M4/M7 & Yüksek performans, FPU, DSP\\
\rowcolor{acikGri} STM32H7 & Cortex-M7 & En yüksek performans (480 MHz+)\\
\bottomrule
\end{tabularx}
\end{center}

\gsub{Neden Register Seviyesi Öğrenmeli?}

Çoğu proje ST'nin HAL (Hardware Abstraction Layer) kütüphanesiyle yazılır ---
hızlıdır ama neyin nasıl çalıştığını gizler. Bu rehber önce \emph{register
seviyesini} (bare-metal) öğretir, çünkü:

\begin{itemize}\itemsep2pt
\item Donanımın gerçekte nasıl çalıştığını anlarsın (datasheet okuma becerisi).
\item HAL'ın gizlediği hataları ayıklayabilirsin.
\item Kritik yerlerde daha hızlı, daha küçük kod yazabilirsin.
\item Assembly'ye (Kısım 5-6) geçiş doğal olur.
\end{itemize}

\bilgi{Öğrenmek için gerçek bir karta ihtiyacın var. Ucuz ve yaygın seçenekler:
\textbf{STM32F103 ``Blue Pill''} (Cortex-M3, çok ucuz), \textbf{STM32 Nucleo}
kartları (dahili ST-Link programlayıcı ile) ve \textbf{STM32 Discovery}
kartları. Programlamak için bir \code{ST-Link} programlayıcı/hata ayıklayıcı
gerekir (Nucleo'da dahili).}

% =====================================================================
\gsec{Bellek Haritası ve Register Modeli}
% =====================================================================

STM32'yi register seviyesinde programlamak, aslında \emph{belirli bellek
adreslerine} değer yazıp okumaktır. Her çevre birimi, bellek haritasında
sabit bir adrese ``eşlenmiştir'' (memory-mapped I/O). Bir GPIO pinini yakmak,
o pine karşılık gelen bir register'a bir bit yazmaktan ibarettir.

\gsub{Cortex-M Bellek Haritası}

Cortex-M, 4 GB'lık düz (flat) bir adres uzayına sahiptir; bu uzay sabit
bölgelere ayrılmıştır:

\begin{center}
\begin{tabularx}{\textwidth}{@{}>{\ttfamily\bfseries\color{anaMavi}}l >{\bfseries}l X@{}}
\toprule
\rowcolor{acikMavi}\color{anaMavi}\textbf{Adres} & \color{anaMavi}\textbf{Bölge} & \color{anaMavi}\textbf{İçerik}\\
\midrule
0x0000\_0000 & Code & Flash (program kodu), vektör tablosu\\
\rowcolor{acikGri} 0x2000\_0000 & SRAM & Değişkenler, yığın (stack), heap\\
0x4000\_0000 & Peripheral & Çevre birimi register'ları (GPIO, TIM\dots)\\
\rowcolor{acikGri} 0xE000\_0000 & Cortex-M iç & NVIC, SysTick, SCB (sistem register'ları)\\
\bottomrule
\end{tabularx}
\end{center}

\gsub{Register'a Erişim: volatile İşaretçiler}

Bir çevre birimi register'ına C'de erişmek için, adresini bir \code{volatile}
işaretçi olarak ele alırız. \code{volatile} anahtar kelimesi kritiktir:
derleyiciye ``bu değer her an donanım tarafından değişebilir, okumaları/yazmaları
optimize etme'' der.

\begin{lstlisting}
#include <stdint.h>

/* GPIOC çevre biriminin taban adresi (STM32F1 için) */
#define GPIOC_BASE   0x40011000UL
#define RCC_BASE     0x40021000UL

/* Register'lara volatile işaretçilerle eriş */
#define RCC_APB2ENR  (*(volatile uint32_t *)(RCC_BASE + 0x18))
#define GPIOC_CRH    (*(volatile uint32_t *)(GPIOC_BASE + 0x04))
#define GPIOC_ODR    (*(volatile uint32_t *)(GPIOC_BASE + 0x0C))

/* Kullanım: bir register'a doğrudan yaz/oku */
/* RCC_APB2ENR |= (1 << 4);   // GPIOC saatini aç (birazdan açıklanacak) */
\end{lstlisting}

\uyari{\code{volatile}'i unutmak, gömülü sistemlerdeki en sinsi hatalardan
biridir. Derleyici, ``değişmeyen'' sandığı bir register okumasını döngüden
çıkarabilir veya bir yazmayı tamamen atabilir --- kod ``bazen çalışır bazen
çalışmaz''. Donanım register'larına erişim \textbf{her zaman} \code{volatile}
olmalıdır.}

\gsub{CMSIS: Standart Register Tanımları}

Yukarıdaki gibi adresleri elle tanımlamak yerine, ARM'ın \emph{CMSIS}
(Cortex Microcontroller Software Interface Standard) standardı ve ST'nin aygıt
başlıkları (\code{stm32f1xx.h}) her register'ı hazır, tip güvenli
\code{struct}'larla tanımlar. Gerçek projede bunları kullanırız.

\begin{lstlisting}
#include "stm32f1xx.h"      /* CMSIS aygıt başlığı: tüm register'lar tanımlı */

/* CMSIS ile aynı işlem çok daha okunaklı: */
/* Her çevre birimi bir struct işaretçisidir (GPIOC, RCC, TIM2...) */
RCC->APB2ENR |= RCC_APB2ENR_IOPCEN;   /* GPIOC saatini aç */
GPIOC->ODR   |= (1 << 13);            /* PC13 pinini yüksek yap */
\end{lstlisting}

\bilgi{CMSIS başlıkları, bir çevre biriminin taban adresini bir \code{struct}
işaretçisi olarak tanımlar; struct alanları da register'lara karşılık gelir.
Örneğin \code{GPIOC->ODR}, derleyici tarafından tam olarak
\code{*(volatile uint32\_t*)(0x40011000 + 0x0C)} adresine çevrilir --- yani
elle yaptığımızla aynı, ama okunaklı ve hatasız.}

% =====================================================================
\gsec{Saat Sistemi ve Araç Zinciri}
% =====================================================================

\gsub{Saat (Clock) Sistemi ve RCC}

Güç tasarrufu için STM32'de çevre birimlerinin saatleri \emph{varsayılan olarak
kapalıdır}. Bir çevre birimini (GPIO, timer, UART\dots) kullanmadan önce, onun
saatini \emph{RCC} (Reset and Clock Control) üzerinden \textbf{açmalısın}. Bu,
yeni başlayanların en sık unuttuğu adımdır.

\begin{lstlisting}
#include "stm32f1xx.h"

/* Bir çevre birimini kullanmadan ÖNCE saatini aç: */
RCC->APB2ENR |= RCC_APB2ENR_IOPCEN;   /* GPIOC (APB2 veri yolunda) */
RCC->APB1ENR |= RCC_APB1ENR_TIM2EN;   /* TIM2 (APB1 veri yolunda) */
RCC->APB2ENR |= RCC_APB2ENR_USART1EN; /* USART1 */
\end{lstlisting}

\uyari{\textbf{``Kodum çalışmıyor'' vakalarının bir numaralı sebebi:} çevre
biriminin saatini açmayı unutmak. Saat kapalıysa register'lara yazdığın değerler
\emph{hiçbir etki yapmaz} (ve okumalar 0 döner). Yeni bir çevre birimi
kullanırken ilk yapman gereken, doğru RCC register'ında saatini açmaktır ---
hangi veri yolunda (AHB/APB1/APB2) olduğunu datasheet'ten kontrol et.}

\gsub{Araç Zinciri ve Derleme}

STM32 için kod, \code{arm-none-eabi-gcc} çapraz derleyicisiyle (cross-compiler)
derlenir --- kodu PC'de derleyip MCU için makine kodu üretir.

\begin{lstlisting}[style=shstyle]
# ARM GCC araç zincirini kur (Debian/Ubuntu)
$ sudo apt install gcc-arm-none-eabi gdb-multiarch openocd stlink-tools

# Derle: startup + linker script gerekir
$ arm-none-eabi-gcc -mcpu=cortex-m3 -mthumb -Wall -g \
    -T stm32f103.ld -nostartfiles \
    startup.s main.c -o firmware.elf

# ELF'ten ikili (binary) üret
$ arm-none-eabi-objcopy -O binary firmware.elf firmware.bin

# Karta yükle (ST-Link ile)
$ st-flash write firmware.bin 0x08000000

# Boyutu gör (Flash/RAM kullanımı)
$ arm-none-eabi-size firmware.elf
\end{lstlisting}

\bilgi{\textbf{Bare-metal derleme üç özel parça gerektirir:} (1) \emph{startup
kodu} (\code{startup.s}) --- reset sonrası çalışan, yığını kuran ve
\code{main}'e atlayan assembly; (2) \emph{linker script} (\code{.ld}) --- kodun
Flash'ta, değişkenlerin RAM'de nereye yerleşeceğini tarif eder; (3)
\emph{vektör tablosu} --- kesme işleyici adreslerini tutar. STM32CubeIDE bunları
otomatik üretir; ama nasıl çalıştıklarını bilmek şarttır. Startup kodunu
Kısım 6'da assembly ile inceleyeceğiz.}

% =====================================================================
\gsec{Saat Ağacı (Clock Tree) Detayı}
% =====================================================================

STM32'nin ne kadar hızlı çalışacağını ve her çevre biriminin hangi frekansı
alacağını \emph{saat ağacı} belirler. Reset sonrası MCU yavaş dahili bir
osilatörle çalışır; tam hıza ulaşmak için saat ağacını yapılandırmak gerekir.

\gsub{Saat Kaynakları ve PLL}

\begin{center}
\begin{tabularx}{\textwidth}{@{}>{\ttfamily\bfseries\color{anaMavi}}l >{\bfseries}l X@{}}
\toprule
\rowcolor{acikMavi}\color{anaMavi}\textbf{Kaynak} & \color{anaMavi}\textbf{Ad} & \color{anaMavi}\textbf{Özellik}\\
\midrule
HSI & Dahili yüksek hız & RC osilatör (ör. 8/16 MHz), hazır ama az hassas.\\
\rowcolor{acikGri} HSE & Harici yüksek hız & Kristal (ör. 8 MHz), hassas.\\
PLL & Faz kilitli çevrim & Kaynağı çarparak yüksek frekans üretir (ör. 72 MHz).\\
\rowcolor{acikGri} LSE / LSI & Düşük hız & 32.768 kHz --- RTC ve watchdog için.\\
\bottomrule
\end{tabularx}
\end{center}

Tipik yol: HSE (8 MHz kristal) $\to$ PLL ile $\times 9$ $\to$ 72 MHz sistem
saati (SYSCLK). Ardından \emph{prescaler}'lar ile AHB, APB1, APB2 veri yolları
için türev saatler üretilir.

\begin{lstlisting}
#include "stm32f1xx.h"

/* HSE + PLL ile 72 MHz'e çık (STM32F103) */
void clock_init_72mhz(void)
{
    RCC->CR |= RCC_CR_HSEON;                    /* harici kristali başlat */
    while (!(RCC->CR & RCC_CR_HSERDY));         /* hazır olmasını bekle */

    FLASH->ACR |= FLASH_ACR_LATENCY_2;         /* 72 MHz için 2 flash bekleme */

    /* PLL: kaynak HSE, çarpan x9 -> 8 MHz * 9 = 72 MHz */
    RCC->CFGR |= RCC_CFGR_PLLSRC | RCC_CFGR_PLLMULL9;
    RCC->CFGR |= RCC_CFGR_PPRE1_DIV2;          /* APB1 max 36 MHz -> /2 */

    RCC->CR |= RCC_CR_PLLON;                    /* PLL'i başlat */
    while (!(RCC->CR & RCC_CR_PLLRDY));

    RCC->CFGR |= RCC_CFGR_SW_PLL;              /* sistem saatini PLL yap */
    while ((RCC->CFGR & RCC_CFGR_SWS) != RCC_CFGR_SWS_PLL);
}
\end{lstlisting}

\uyari{Saat frekansını artırırken \textbf{Flash bekleme durumunu} (wait state,
\code{FLASH\_ACR\_LATENCY}) doğru ayarlamalısın: Flash belleği CPU kadar hızlı
okunamaz, bu yüzden yüksek frekansta ek bekleme gerekir. Ayarı unutmak, MCU'nun
kilitlenmesine veya kararsız çalışmasına yol açar. Her frekans aralığı için
gereken bekleme sayısı datasheet'te belirtilir.}

\bilgi{Saat ağacını elle kurmak yorucudur; ST'nin \textbf{STM32CubeMX} aracı,
grafik bir saat ağacı editörüyle bu kodu otomatik üretir. Yine de mekanizmayı
bilmek, ``çevre birimim yanlış frekansta çalışıyor'' tarzı hataları çözmek için
gereklidir. Timer ve UART frekansları doğrudan APB saatlerine bağlıdır.}

% #####################################################################
\gpart{2}{GPIO ve Register Seviyesi Programlama}
% #####################################################################

% =====================================================================
\gsec{GPIO: Genel Amaçlı Giriş/Çıkış}
% =====================================================================

GPIO (General-Purpose Input/Output) pinleri, MCU'nun dış dünyayla en temel
bağlantısıdır: bir LED yakmak, bir butonu okumak, bir röleyi tetiklemek hep
GPIO ile yapılır. Her pin ya \emph{çıkış} (output --- MCU pini sürer) ya da
\emph{giriş} (input --- MCU pini okur) olarak yapılandırılır.

\gsub{GPIO Register'ları}

Her GPIO portu (GPIOA, GPIOB\dots) birkaç register'a sahiptir. En önemlileri:

\begin{center}
\begin{tabularx}{\textwidth}{@{}>{\ttfamily\bfseries\color{anaMavi}}l X@{}}
\toprule
\rowcolor{acikMavi}\color{anaMavi}\textbf{Register} & \color{anaMavi}\textbf{İşlevi}\\
\midrule
CRL / CRH & Yapılandırma: pin giriş mi çıkış mı, hız, mod (F1 serisi).\\
\rowcolor{acikGri} MODER & Pin modu (F0/F4 serisinde CRL/CRH yerine).\\
IDR & Input Data Register: giriş pinlerinin durumunu \emph{okur}.\\
\rowcolor{acikGri} ODR & Output Data Register: çıkış pinlerine değer \emph{yazar}.\\
BSRR & Bit Set/Reset: pinleri atomik olarak 1 veya 0 yapar.\\
\bottomrule
\end{tabularx}
\end{center}

\gsub{Bit İşlemleri: Register Programlamanın Dili}

Register'larla çalışmak, tek tek \emph{bitleri} ayarlamaktır. Diğer bitleri
bozmadan belirli bitleri değiştirmek için bit maskeleri kullanılır. Bu dört
işlem, gömülü programlamanın temelidir:

\begin{center}
\begin{tabularx}{\textwidth}{@{}>{\bfseries\color{anaMavi}}l >{\ttfamily}l X@{}}
\toprule
\rowcolor{acikMavi}\color{anaMavi}\textbf{İşlem} & \color{anaMavi}\textbf{Kod} & \color{anaMavi}\textbf{Anlamı}\\
\midrule
Bit kur (set) & REG |= (1 << n); & $n$. biti 1 yap (diğerlerine dokunma).\\
\rowcolor{acikGri} Bit temizle (clear) & REG \&= \~{}(1 << n); & $n$. biti 0 yap.\\
Bit çevir (toggle) & REG \^{}= (1 << n); & $n$. biti tersle.\\
\rowcolor{acikGri} Bit oku (test) & (REG >> n) \& 1 & $n$. bit 1 mi?\\
\bottomrule
\end{tabularx}
\end{center}

\ipucu{\textbf{``Read-modify-write''} kalıbı: \code{REG |= (1 << n)} aslında
üç işlemdir --- register'ı oku, biti değiştir, geri yaz. Bu yüzden bir kesme
araya girerse yarış koşulu olabilir. STM32'nin \code{BSRR} register'ı bunu
çözer: tek bir yazmayla, atomik olarak pin set/reset yapar --- read-modify-write
gerekmez.}

% =====================================================================
\gsec{İlk Program: LED Yakma (Blink)}
% =====================================================================

Gömülü dünyanın ``merhaba dünya''sı bir LED yakıp söndürmektir. STM32F103
``Blue Pill'' kartında dahili LED \emph{PC13} pinine bağlıdır. Adım adım
register seviyesinde yapalım.

\gsub{Bare-Metal LED Blink}

\begin{lstlisting}
#include "stm32f1xx.h"

/* Basit gecikme: işlemciyi meşgul tutan bir döngü (kaba, kalibre değil) */
static void delay(volatile uint32_t count)
{
    while (count--) {
        __asm__("nop");   /* "no operation" -- derleyici döngüyü silmesin */
    }
}

int main(void)
{
    /* 1) GPIOC saatini aç (RCC üzerinden) */
    RCC->APB2ENR |= RCC_APB2ENR_IOPCEN;

    /* 2) PC13'ü çıkış olarak yapılandır (F1: CRH, pin 13 -> bit 20-23)
          MODE = 0b10 (2 MHz çıkış), CNF = 0b00 (push-pull) */
    GPIOC->CRH &= ~(0xF << 20);        /* önce PC13 alanını temizle */
    GPIOC->CRH |=  (0x2 << 20);        /* sonra çıkış modunu kur */

    /* 3) Sonsuz döngüde LED'i çevir */
    while (1) {
        GPIOC->ODR ^= (1 << 13);       /* PC13'ü tersle (toggle) */
        delay(500000);                 /* bir süre bekle */
    }
}
\end{lstlisting}

\notk[BSRR ile atomik kontrol]{Yukarıda \code{ODR \^{}= ...} kullandık; ama
pini kesin olarak açıp kapamak için \code{BSRR} daha güvenlidir. BSRR'nin alt
16 biti pini \emph{set} eder (1 yapar), üst 16 biti \emph{reset} eder (0 yapar):
\code{GPIOC->BSRR = (1 << 13);} pini yakar,
\code{GPIOC->BSRR = (1 << (13+16));} söndürür --- her ikisi de tek, atomik
yazmadır.}

\gsub{Aynı İş HAL ile}

Karşılaştırma için, aynı işi ST'nin HAL kütüphanesi çok daha yüksek seviyede
yapar (ama arkada aynı register'lara yazar):

\begin{lstlisting}
#include "stm32f1xx_hal.h"

int main(void)
{
    HAL_Init();
    __HAL_RCC_GPIOC_CLK_ENABLE();          /* saati aç */

    GPIO_InitTypeDef cfg = {0};
    cfg.Pin   = GPIO_PIN_13;
    cfg.Mode  = GPIO_MODE_OUTPUT_PP;       /* push-pull çıkış */
    cfg.Speed = GPIO_SPEED_FREQ_LOW;
    HAL_GPIO_Init(GPIOC, &cfg);

    while (1) {
        HAL_GPIO_TogglePin(GPIOC, GPIO_PIN_13);
        HAL_Delay(500);                    /* 500 ms (SysTick tabanlı) */
    }
}
\end{lstlisting}

\bilgi{\textbf{Bare-metal vs HAL:} Bare-metal kod küçük, hızlı ve şeffaftır ama
her register'ı elle ayarlaman gerekir (datasheet okuma). HAL taşınabilir ve
hızlı geliştirme sağlar ama daha büyük/yavaştır ve bazen davranışı gizler.
Öğrenirken bare-metal, üretimde çoğu zaman HAL (veya LL --- Low-Layer)
kullanılır. İkisini de anlamak en iyisidir.}

% =====================================================================
\gsec{Giriş Okuma: Buton}
% =====================================================================

Bir butonu okumak, bir GPIO pinini \emph{giriş} olarak yapılandırıp \code{IDR}
register'ından durumunu okumaktır. Butonun elektriksel olarak tanımlı bir
seviyede durması için \emph{pull-up} veya \emph{pull-down} direnci gerekir
(STM32'de dahili olarak etkinleştirilebilir).

\begin{lstlisting}
#include "stm32f1xx.h"

int main(void)
{
    /* GPIOA ve GPIOC saatlerini aç */
    RCC->APB2ENR |= RCC_APB2ENR_IOPAEN | RCC_APB2ENR_IOPCEN;

    /* PA0'ı giriş + pull-up olarak yapılandır (CRL, pin 0 -> bit 0-3)
       MODE=0b00 (giriş), CNF=0b10 (pull-up/pull-down) */
    GPIOA->CRL &= ~(0xF << 0);
    GPIOA->CRL |=  (0x8 << 0);
    GPIOA->ODR |=  (1 << 0);            /* ODR biti pull-UP seçer */

    /* PC13'ü çıkış yap (LED) */
    GPIOC->CRH &= ~(0xF << 20);
    GPIOC->CRH |=  (0x2 << 20);

    while (1) {
        /* Butona basılınca PA0 düşer (pull-up olduğu için basılı = 0) */
        if (((GPIOA->IDR >> 0) & 1) == 0) {
            GPIOC->BSRR = (1 << 13);       /* LED aç */
        } else {
            GPIOC->BSRR = (1 << (13 + 16)); /* LED kapat */
        }
    }
}
\end{lstlisting}

\uyari{\textbf{Buton sıçraması (debouncing):} Mekanik butonlar, basıldığında
birkaç milisaniye boyunca hızla açılıp kapanır (bounce). Ham okuma tek bir
basışı ``çok basış'' olarak algılayabilir. Çözüm: bir okuma sonrası kısa bir
gecikme (ör. 20 ms) koymak veya durumu birkaç kez üst üste teyit etmek
(software debouncing).}

% #####################################################################
\gpart{3}{Çevre Birimleri (Peripherals)}
% #####################################################################

STM32'nin gücü, çip içindeki zengin çevre birimlerinden gelir: zamanlayıcılar,
seri iletişim (UART, SPI, I2C), analog-dijital çevirici (ADC) ve daha fazlası.
Her biri kendi register'larıyla yapılandırılır. Bu kısım en yaygın olanları
tanıtır.

% =====================================================================
\gsec{Zamanlayıcılar (Timers) ve PWM}
% =====================================================================

Zamanlayıcı (timer), belirli bir saat frekansıyla sayan bir sayaçtır. Kesin
zaman aralıkları üretmek, olayları saymak, gecikme oluşturmak ve \emph{PWM}
(darbe genişlik modülasyonu) üretmek için kullanılır. PWM, LED parlaklığı
ayarlama ve motor/servo kontrolünün temelidir.

\gsub{Zamanlayıcı Nasıl Çalışır?}

Bir timer, \code{PSC} (prescaler --- ön bölücü) ile saati böler, sonra
\code{ARR} (auto-reload register) değerine kadar sayar; oraya ulaşınca sıfırlanıp
bir olay (ve isteğe bağlı kesme) üretir. Timer frekansı:
$f_{\text{timer}} = \dfrac{f_{\text{clock}}}{(PSC+1)\times(ARR+1)}$.

\begin{lstlisting}
#include "stm32f1xx.h"

/* TIM2'yi 1 Hz periyotla (72 MHz saat varsayımı) yapılandır */
void timer_init(void)
{
    RCC->APB1ENR |= RCC_APB1ENR_TIM2EN;    /* TIM2 saatini aç */

    TIM2->PSC = 7200 - 1;                  /* 72 MHz / 7200 = 10 kHz */
    TIM2->ARR = 10000 - 1;                 /* 10 kHz / 10000 = 1 Hz */
    TIM2->CR1 |= TIM_CR1_CEN;              /* sayacı başlat (counter enable) */
}

/* Sayacın taşmasını (update event) yoklayarak bekle */
void timer_wait_update(void)
{
    while (!(TIM2->SR & TIM_SR_UIF));      /* update flag'i bekle */
    TIM2->SR &= ~TIM_SR_UIF;               /* flag'i temizle */
}
\end{lstlisting}

\gsub{PWM Üretme}

PWM, sabit frekanslı ama değişken \emph{doluluk oranına} (duty cycle) sahip bir
kare dalgadır. \code{CCR} (capture/compare register) doluluk oranını belirler.
Örneğin bir LED'in parlaklığını, göz ortalamayı algıladığı için, PWM ile
ayarlarız.

\begin{lstlisting}
/* TIM3 Kanal 1'de PWM (PA6 pini) -- LED parlaklık kontrolü */
void pwm_init(void)
{
    RCC->APB1ENR |= RCC_APB1ENR_TIM3EN;
    RCC->APB2ENR |= RCC_APB2ENR_IOPAEN;

    /* PA6'yı alternatif fonksiyon push-pull çıkış yap (CRL bit 24-27) */
    GPIOA->CRL &= ~(0xF << 24);
    GPIOA->CRL |=  (0xB << 24);            /* MODE=0b11 (50MHz), CNF=0b10 (AF-PP) */

    TIM3->PSC = 72 - 1;                    /* 72 MHz / 72 = 1 MHz sayım */
    TIM3->ARR = 1000 - 1;                  /* 1 MHz / 1000 = 1 kHz PWM */
    TIM3->CCR1 = 250;                      /* doluluk = 250/1000 = %25 */
    TIM3->CCMR1 |= (6 << 4);               /* PWM mode 1 (OC1M = 110) */
    TIM3->CCER  |= TIM_CCER_CC1E;          /* kanal 1 çıkışını etkinleştir */
    TIM3->CR1   |= TIM_CR1_CEN;
}
/* Parlaklığı değiştirmek için: TIM3->CCR1 = yeni_deger; (0..999) */
\end{lstlisting}

\ipucu{PWM'in doluluk oranını (\code{CCR}) çalışırken değiştirerek bir LED'i
yavaşça parlatıp söndürebilir (fade), bir servo motorun açısını ayarlayabilir
veya bir DC motorun hızını kontrol edebilirsin. Frekans (\code{PSC}/\code{ARR})
sabit kalır, yalnızca \code{CCR} değişir.}

\gsub{HAL Driver ile Timer ve PWM}

\gsubsub{Periyodik Timer Kesmesi}

HAL ile bir timer'ı periyodik kesme üretecek şekilde kurmak birkaç satırdır;
her taşmada \code{HAL\_TIM\_PeriodElapsedCallback} çağrılır.

\begin{lstlisting}
#include "stm32f1xx_hal.h"

TIM_HandleTypeDef htim2;

void timer_hal_init(void)
{
    __HAL_RCC_TIM2_CLK_ENABLE();

    htim2.Instance           = TIM2;
    htim2.Init.Prescaler     = 7200 - 1;   /* 72 MHz / 7200 = 10 kHz */
    htim2.Init.Period        = 10000 - 1;  /* 10 kHz / 10000 = 1 Hz */
    htim2.Init.CounterMode   = TIM_COUNTERMODE_UP;
    HAL_TIM_Base_Init(&htim2);

    HAL_TIM_Base_Start_IT(&htim2);         /* kesmeli olarak başlat */
    HAL_NVIC_EnableIRQ(TIM2_IRQn);
}

/* Her periyot dolduğunda (1 Hz) çağrılır */
void HAL_TIM_PeriodElapsedCallback(TIM_HandleTypeDef *htim)
{
    if (htim->Instance == TIM2) {
        HAL_GPIO_TogglePin(GPIOC, GPIO_PIN_13);   /* LED çevir */
    }
}

void TIM2_IRQHandler(void) { HAL_TIM_IRQHandler(&htim2); }
\end{lstlisting}

\gsubsub{HAL ile PWM}

\begin{lstlisting}
TIM_HandleTypeDef htim3;

void pwm_hal_init(void)
{
    __HAL_RCC_TIM3_CLK_ENABLE();

    htim3.Instance        = TIM3;
    htim3.Init.Prescaler  = 72 - 1;        /* 1 MHz sayım */
    htim3.Init.Period     = 1000 - 1;      /* 1 kHz PWM */
    HAL_TIM_PWM_Init(&htim3);

    TIM_OC_InitTypeDef oc = {0};
    oc.OCMode     = TIM_OCMODE_PWM1;
    oc.Pulse      = 250;                    /* doluluk = %25 (250/1000) */
    oc.OCPolarity = TIM_OCPOLARITY_HIGH;
    HAL_TIM_PWM_ConfigChannel(&htim3, &oc, TIM_CHANNEL_1);

    HAL_TIM_PWM_Start(&htim3, TIM_CHANNEL_1);   /* PWM çıkışını başlat */
}

/* Çalışırken doluluk oranını değiştir (ör. LED parlaklığı / motor hızı) */
void pwm_set_duty(uint16_t duty)           /* 0..999 */
{
    __HAL_TIM_SET_COMPARE(&htim3, TIM_CHANNEL_1, duty);
}
\end{lstlisting}

\gsubsub{Input Capture: Sinyal Ölçme}

Timer'ın bir başka gücü \emph{input capture}'dır: bir giriş pinindeki kenar
geldiğinde sayaç değerini yakalar. Böylece bir sinyalin frekansını veya darbe
genişliğini (ör. HC-SR04 mesafe sensörü, RC alıcı) ölçebilirsin.

\begin{lstlisting}
/* Kenar yakalandığında çağrılır -- iki yakalama arası fark periyodu verir */
void HAL_TIM_IC_CaptureCallback(TIM_HandleTypeDef *htim)
{
    if (htim->Channel == HAL_TIM_ACTIVE_CHANNEL_1) {
        uint32_t captured = HAL_TIM_ReadCapturedValue(htim, TIM_CHANNEL_1);
        /* Ardışık iki 'captured' farkı -> periyot -> frekans hesaplanır */
    }
}
/* Başlatma: HAL_TIM_IC_Start_IT(&htim, TIM_CHANNEL_1); */
\end{lstlisting}

\notk[Encoder modu]{Timer'lar ayrıca \emph{encoder modunda} çalışabilir:
bir rotary encoder'ın (kadran/motor mili) iki fazını (A, B) sayarak konum ve
yön bilgisini donanımda üretir. HAL'da \code{HAL\_TIM\_Encoder\_Start()} ile
başlatılır ve sayaç değeri (\code{\_\_HAL\_TIM\_GET\_COUNTER}) doğrudan pozisyonu
verir --- CPU hiç yorulmaz.}

% =====================================================================
\gsec{Seri İletişim: UART}
% =====================================================================

UART (Universal Asynchronous Receiver/Transmitter), iki cihaz arasında iki tel
(TX, RX) üzerinden seri veri gönderen en yaygın iletişim yöntemidir. PC ile
hata ayıklama mesajı alışverişinde (\code{printf} tarzı) vazgeçilmezdir.

\gsub{UART Yapılandırma ve Gönderme}

UART'ın en önemli ayarı \emph{baud rate}'tir (saniyedeki bit sayısı, ör. 9600
veya 115200) --- iki taraf aynı olmalıdır. Baud, \code{BRR} register'ıyla
ayarlanır.

\begin{lstlisting}
#include "stm32f1xx.h"

void uart_init(void)
{
    RCC->APB2ENR |= RCC_APB2ENR_USART1EN | RCC_APB2ENR_IOPAEN;

    /* PA9 = TX (alternatif fonksiyon push-pull) */
    GPIOA->CRH &= ~(0xF << 4);
    GPIOA->CRH |=  (0xB << 4);

    /* Baud = 115200 @ 72 MHz -> BRR = 72e6 / 115200 = 625 */
    USART1->BRR = 625;
    USART1->CR1 |= USART_CR1_TE;           /* verici (transmitter) etkin */
    USART1->CR1 |= USART_CR1_UE;           /* USART'ı etkinleştir */
}

/* Tek bir karakter gönder (gönderme tamponu boşalana kadar bekle) */
void uart_send_char(char c)
{
    while (!(USART1->SR & USART_SR_TXE));  /* TXE: tampon boş mu? */
    USART1->DR = c;                        /* veriyi yaz -> otomatik gönderilir */
}

/* Bir string gönder */
void uart_send_string(const char *s)
{
    while (*s) {
        uart_send_char(*s++);
    }
}
\end{lstlisting}

\bilgi{UART ile PC'ye mesaj göndermek, gömülü hata ayıklamanın en pratik
yoludur (bir tür ``\code{printf} debugging''). Bir USB-Serial dönüştürücü
(FTDI/CP2102) ile MCU'nun TX pinini PC'ye bağlar, PC'de bir terminal
(\code{minicom}, PuTTY, \code{screen}) açarsın. \code{printf}'i UART'a
yönlendirmek için \code{\_write} sistem çağrısını yeniden tanımlayabilirsin.}

\gsub{HAL Driver ile UART}

Gerçek projelerde UART'ı register seviyesinde yönetmek yerine ST'nin
\emph{HAL} (Hardware Abstraction Layer) sürücüsü kullanılır. HAL, tüm register
detaylarını gizler; sadece bir \emph{handle} yapısını (\code{UART\_HandleTypeDef})
doldurup hazır fonksiyonları çağırırsın. Aşağıda üç çalışma modu gösterilir:
\emph{blocking} (yoklamalı), \emph{interrupt} ve \emph{DMA}.

\gsubsub{Yapılandırma (Init)}

\begin{lstlisting}
#include "stm32f1xx_hal.h"

UART_HandleTypeDef huart1;                 /* UART handle (durum + ayarlar) */

void uart_hal_init(void)
{
    __HAL_RCC_USART1_CLK_ENABLE();         /* saatleri aç */
    __HAL_RCC_GPIOA_CLK_ENABLE();

    /* PA9=TX, PA10=RX pinlerini alternatif fonksiyon olarak ayarla */
    GPIO_InitTypeDef gpio = {0};
    gpio.Pin   = GPIO_PIN_9;
    gpio.Mode  = GPIO_MODE_AF_PP;          /* alternatif fonksiyon push-pull */
    gpio.Speed = GPIO_SPEED_FREQ_HIGH;
    HAL_GPIO_Init(GPIOA, &gpio);
    gpio.Pin  = GPIO_PIN_10;
    gpio.Mode = GPIO_MODE_INPUT;           /* RX giriş */
    HAL_GPIO_Init(GPIOA, &gpio);

    /* UART parametreleri */
    huart1.Instance          = USART1;
    huart1.Init.BaudRate     = 115200;     /* baud hızı */
    huart1.Init.WordLength   = UART_WORDLENGTH_8B;
    huart1.Init.StopBits     = UART_STOPBITS_1;
    huart1.Init.Parity       = UART_PARITY_NONE;
    huart1.Init.Mode         = UART_MODE_TX_RX;   /* hem gönder hem al */
    huart1.Init.HwFlowCtl    = UART_HWCONTROL_NONE;
    HAL_UART_Init(&huart1);                /* tüm register'ları HAL ayarlar */
}
\end{lstlisting}

\gsubsub{Blocking (Yoklamalı) Gönderme ve Alma}

En basit mod: fonksiyon, işlem bitene (veya zaman aşımına) kadar \emph{bekler}.

\begin{lstlisting}
uint8_t rx_byte;
char msg[] = "Merhaba UART\r\n";

/* Gönder: işlem bitene kadar bekle (son parametre = zaman aşımı ms) */
HAL_UART_Transmit(&huart1, (uint8_t *)msg, sizeof(msg) - 1, HAL_MAX_DELAY);

/* Al: 1 bayt gelene kadar bekle */
HAL_UART_Receive(&huart1, &rx_byte, 1, HAL_MAX_DELAY);
\end{lstlisting}

\gsubsub{Interrupt Modu (Non-Blocking)}

Interrupt modunda CPU beklemez; veri gelince/gidince HAL bir \emph{callback}
çağırır. Bu, işlemciyi bloklamadan arka planda iletişim sağlar.

\begin{lstlisting}
uint8_t rx_buffer[1];

void uart_start_receive(void)
{
    /* 1 bayt için kesmeli alım başlat -- hemen döner, beklemez */
    HAL_UART_Receive_IT(&huart1, rx_buffer, 1);
}

/* Alım tamamlanınca HAL bunu otomatik çağırır (callback'i sen yazarsın) */
void HAL_UART_RxCpltCallback(UART_HandleTypeDef *huart)
{
    if (huart->Instance == USART1) {
        /* Gelen baytı işle (örneğin geri gönder - echo) */
        HAL_UART_Transmit_IT(&huart1, rx_buffer, 1);
        HAL_UART_Receive_IT(&huart1, rx_buffer, 1);  /* tekrar dinle */
    }
}

/* USART1 kesmesi HAL'ın işleyicisine yönlendirilmeli (ISR): */
void USART1_IRQHandler(void)
{
    HAL_UART_IRQHandler(&huart1);          /* HAL kesmeyi çözer, callback'i çağırır */
}
\end{lstlisting}

\gsubsub{DMA Modu (Yüksek Verim)}

Büyük veri bloklarını CPU'yu hiç meşgul etmeden aktarmak için DMA kullanılır.

\begin{lstlisting}
/* Bir bloğu DMA ile gönder -- CPU başka iş yaparken transfer arka planda olur */
HAL_UART_Transmit_DMA(&huart1, (uint8_t *)msg, sizeof(msg) - 1);

/* DMA transferi bitince çağrılır */
void HAL_UART_TxCpltCallback(UART_HandleTypeDef *huart)
{
    /* Gönderim tamamlandı -- bir sonraki bloğu başlatabilirsin */
}
\end{lstlisting}

\gsubsub{printf'i UART'a Yönlendirme}

\code{printf}'in UART'a çıkmasını sağlamak için \code{\_write} (veya
\code{\_\_io\_putchar}) sistem fonksiyonunu HAL üzerinden yeniden tanımlarız:

\begin{lstlisting}
#include <stdio.h>

/* newlib'in printf'i sonunda buraya gelir -> UART'a yaz */
int _write(int file, char *ptr, int len)
{
    HAL_UART_Transmit(&huart1, (uint8_t *)ptr, len, HAL_MAX_DELAY);
    return len;
}
/* Artık: printf("Sicaklik: %d C\n", temp);  -> UART terminaline çıkar */
\end{lstlisting}

\bilgi{\textbf{Hangi mod ne zaman?} \emph{Blocking} en basittir, kısa mesajlar
için idealdir ama CPU'yu bekletir. \emph{Interrupt} modu, arka planda sürekli
komut/veri alışverişi (ör. bir konsol) için uygundur. \emph{DMA}, yüksek hızlı
veya büyük bloklar (ör. sensör akışı, log) için CPU'yu tamamen serbest bırakır.
HAL'ın güzelliği, aynı sürücüyle üç modu da kolayca kullanabilmendir.}

% =====================================================================
\gsec{ADC ve Diğer Veri Yolları (SPI, I2C)}
% =====================================================================

\gsub{ADC: Analog Okuma}

ADC (Analog-to-Digital Converter), bir pindeki analog gerilimi (0--3.3 V) bir
sayıya çevirir. Potansiyometre, sıcaklık sensörü, ışık sensörü gibi analog
kaynakları okumak için kullanılır. STM32'nin 12-bit ADC'si 0--4095 arası değer
üretir.

\begin{lstlisting}
#include "stm32f1xx.h"

void adc_init(void)
{
    RCC->APB2ENR |= RCC_APB2ENR_ADC1EN | RCC_APB2ENR_IOPAEN;
    /* PA0 varsayılan olarak analog giriş yapabilir (CRL sıfırla) */
    GPIOA->CRL &= ~(0xF << 0);             /* PA0 = analog mod (0b0000) */

    ADC1->CR2 |= ADC_CR2_ADON;             /* ADC'yi aç */
}

/* Kanal 0'dan (PA0) bir örnek oku */
uint16_t adc_read(void)
{
    ADC1->SQR3 = 0;                        /* dönüşüm sırası: kanal 0 */
    ADC1->CR2 |= ADC_CR2_ADON;             /* dönüşümü başlat */
    while (!(ADC1->SR & ADC_SR_EOC));      /* dönüşüm bitti mi? (EOC) */
    return ADC1->DR;                       /* 12-bit sonucu oku (0..4095) */
}

/* Değeri gerilime çevir: voltage = adc_read() * 3.3 / 4095 */
\end{lstlisting}

\gsub{HAL Driver ile ADC}

Aynı ADC okumasını HAL sürücüsüyle yapmak hem daha okunaklı hem de kesme/DMA
modlarına kolayca geçilebilir. Bir \code{ADC\_HandleTypeDef} handle'ı yapılandırıp
hazır fonksiyonları çağırırız.

\gsubsub{Yapılandırma ve Yoklamalı (Polling) Okuma}

\begin{lstlisting}
#include "stm32f1xx_hal.h"

ADC_HandleTypeDef hadc1;                   /* ADC handle */

void adc_hal_init(void)
{
    __HAL_RCC_ADC1_CLK_ENABLE();
    __HAL_RCC_GPIOA_CLK_ENABLE();

    GPIO_InitTypeDef gpio = {0};
    gpio.Pin  = GPIO_PIN_0;
    gpio.Mode = GPIO_MODE_ANALOG;          /* PA0 = analog giriş */
    HAL_GPIO_Init(GPIOA, &gpio);

    hadc1.Instance                   = ADC1;
    hadc1.Init.ContinuousConvMode    = DISABLE;   /* tek seferlik dönüşüm */
    hadc1.Init.ScanConvMode          = DISABLE;
    hadc1.Init.DataAlign             = ADC_DATAALIGN_RIGHT;
    hadc1.Init.NbrOfConversion       = 1;
    HAL_ADC_Init(&hadc1);

    /* Kanal 0'ı (PA0) dönüşüm sırasına ekle */
    ADC_ChannelConfTypeDef ch = {0};
    ch.Channel      = ADC_CHANNEL_0;
    ch.Rank         = 1;
    ch.SamplingTime = ADC_SAMPLETIME_55CYCLES_5;
    HAL_ADC_ConfigChannel(&hadc1, &ch);
}

/* Yoklamalı okuma: dönüşümü başlat, bitmesini bekle, değeri al */
uint16_t adc_hal_read(void)
{
    HAL_ADC_Start(&hadc1);
    HAL_ADC_PollForConversion(&hadc1, HAL_MAX_DELAY);   /* dönüşümü bekle */
    uint16_t value = HAL_ADC_GetValue(&hadc1);          /* 12-bit sonuç */
    HAL_ADC_Stop(&hadc1);
    return value;
}
\end{lstlisting}

\gsubsub{Interrupt Modu}

\begin{lstlisting}
volatile uint16_t adc_value;

void adc_start_it(void)
{
    HAL_ADC_Start_IT(&hadc1);              /* kesmeli dönüşüm başlat, beklemez */
}

/* Dönüşüm bitince HAL bunu otomatik çağırır */
void HAL_ADC_ConvCpltCallback(ADC_HandleTypeDef *hadc)
{
    adc_value = HAL_ADC_GetValue(hadc);    /* sonucu al */
}

void ADC1_2_IRQHandler(void)
{
    HAL_ADC_IRQHandler(&hadc1);
}
\end{lstlisting}

\gsubsub{DMA ile Sürekli Çok-Kanal Okuma}

Birden çok analog kanalı sürekli ve CPU'suz okumanın en verimli yolu DMA'dır.
HAL, DMA'yı ADC'ye bağlayıp bir tampona otomatik yazar.

\begin{lstlisting}
uint16_t adc_buffer[3];                    /* 3 kanal için tampon */

/* ScanConvMode=ENABLE, ContinuousConvMode=ENABLE, NbrOfConversion=3 ile
   init edilmiş varsayılır; 3 kanal rank 1,2,3'e atanır. */
void adc_start_dma(void)
{
    /* Dönüşümleri başlat; her sonuç otomatik adc_buffer[]'a yazılır */
    HAL_ADC_Start_DMA(&hadc1, (uint32_t *)adc_buffer, 3);
}
/* Artık adc_buffer[0..2] her zaman en güncel 3 kanal değerini tutar -- CPU serbest */
\end{lstlisting}

\ipucu{DMA + \emph{scan mode} kombinasyonu, birden çok sensörü (ör. 3 eksenli
potansiyometre veya çoklu sıcaklık) tek ADC ile sürekli okumanın standart
yoludur. Değerler arka planda tampona akar; ana kod sadece tampondan en güncel
değeri okur. Bu kalıp gömülü sistemlerde çok yaygındır.}

\gsub{SPI ve I2C: Sensör Veri Yolları}

Çoğu sensör ve harici çip (ekran, EEPROM, IMU) MCU ile \emph{SPI} veya
\emph{I2C} veri yolu üzerinden konuşur. İkisi de senkron seri protokoldür ama
farklı ödünleşmeleri vardır.

\begin{center}
\begin{tabularx}{\textwidth}{@{}>{\bfseries\color{anaMavi}}l X X@{}}
\toprule
\rowcolor{acikMavi}\color{anaMavi}\textbf{Özellik} & \color{anaMavi}\textbf{SPI} & \color{anaMavi}\textbf{I2C}\\
\midrule
Tel sayısı & 4 (MOSI, MISO, SCK, CS) & 2 (SDA, SCL)\\
\rowcolor{acikGri} Hız & Çok hızlı (10+ MHz) & Orta (100k--400k--1M)\\
Cihaz sayısı & Her cihaza ayrı CS & Adresle çok cihaz (aynı 2 tel)\\
\rowcolor{acikGri} Karmaşıklık & Basit & Adresleme + ACK protokolü\\
Kullanım & Ekran, SD kart, hızlı ADC & Sensör, EEPROM, RTC\\
\bottomrule
\end{tabularx}
\end{center}

\begin{lstlisting}
/* SPI ile tek bayt gönder/al (full-duplex: aynı anda hem gönderir hem alır) */
uint8_t spi_transfer(uint8_t data)
{
    while (!(SPI1->SR & SPI_SR_TXE));      /* gönderme tamponu boş mu? */
    SPI1->DR = data;                       /* veriyi gönder */
    while (!(SPI1->SR & SPI_SR_RXNE));     /* yanıt geldi mi? */
    return SPI1->DR;                       /* gelen baytı oku */
}
\end{lstlisting}

\gsubsub{HAL ile SPI}

HAL, SPI protokolünün tüm zamanlama detaylarını gizler. Bir
\code{SPI\_HandleTypeDef} yapılandırıp gönder/al fonksiyonlarını çağırırsın.

\begin{lstlisting}
#include "stm32f1xx_hal.h"

SPI_HandleTypeDef hspi1;

void spi_hal_init(void)
{
    __HAL_RCC_SPI1_CLK_ENABLE();
    /* ... GPIO: SCK, MOSI, MISO pinlerini AF olarak ayarla ... */

    hspi1.Instance           = SPI1;
    hspi1.Init.Mode          = SPI_MODE_MASTER;      /* ana (master) cihaz */
    hspi1.Init.Direction     = SPI_DIRECTION_2LINES; /* full-duplex */
    hspi1.Init.DataSize      = SPI_DATASIZE_8BIT;
    hspi1.Init.CLKPolarity   = SPI_POLARITY_LOW;     /* SPI mode 0 */
    hspi1.Init.CLKPhase      = SPI_PHASE_1EDGE;
    hspi1.Init.NSS           = SPI_NSS_SOFT;         /* CS'i yazılım yönetir */
    hspi1.Init.BaudRatePrescaler = SPI_BAUDRATEPRESCALER_8;
    HAL_SPI_Init(&hspi1);
}

void spi_hal_example(void)
{
    uint8_t tx[2] = {0x9F, 0x00};          /* gönderilecek komut */
    uint8_t rx[2];

    /* CS'i indir (bir GPIO pini) -> aktarım -> CS'i kaldır */
    HAL_GPIO_WritePin(GPIOA, GPIO_PIN_4, GPIO_PIN_RESET);
    HAL_SPI_TransmitReceive(&hspi1, tx, rx, 2, HAL_MAX_DELAY);  /* aynı anda gönder+al */
    HAL_GPIO_WritePin(GPIOA, GPIO_PIN_4, GPIO_PIN_SET);
    /* Sadece göndermek için: HAL_SPI_Transmit(&hspi1, tx, 2, HAL_MAX_DELAY); */
}
\end{lstlisting}

\gsubsub{HAL ile I2C}

I2C'de her cihazın bir \emph{adresi} vardır. HAL, adresleme + ACK protokolünü
otomatik yönetir. Özellikle \emph{register okuma/yazma} için \code{Mem\_Read}/
\code{Mem\_Write} fonksiyonları çok pratiktir (sensörlerin çoğu register
tabanlıdır).

\begin{lstlisting}
I2C_HandleTypeDef hi2c1;

void i2c_hal_init(void)
{
    __HAL_RCC_I2C1_CLK_ENABLE();
    /* ... GPIO: SCL, SDA pinlerini AF open-drain olarak ayarla ... */

    hi2c1.Instance       = I2C1;
    hi2c1.Init.ClockSpeed = 100000;        /* 100 kHz standart mod */
    hi2c1.Init.DutyCycle  = I2C_DUTYCYCLE_2;
    hi2c1.Init.OwnAddress1 = 0;
    hi2c1.Init.AddressingMode = I2C_ADDRESSINGMODE_7BIT;
    HAL_I2C_Init(&hi2c1);
}

void i2c_hal_example(void)
{
    uint8_t dev_addr = 0x68 << 1;          /* 7-bit adres, HAL 8-bit ister */
    uint8_t data;

    /* Bir sensör register'ından oku (ör. MPU6050 WHO_AM_I = reg 0x75) */
    HAL_I2C_Mem_Read(&hi2c1, dev_addr, 0x75, I2C_MEMADD_SIZE_8BIT,
                     &data, 1, HAL_MAX_DELAY);

    /* Bir register'a yaz */
    uint8_t config = 0x00;
    HAL_I2C_Mem_Write(&hi2c1, dev_addr, 0x6B, I2C_MEMADD_SIZE_8BIT,
                      &config, 1, HAL_MAX_DELAY);
}
\end{lstlisting}

\bilgi{SPI ve I2C'yi register seviyesinde elle yönetmek zahmetlidir (zamanlama,
protokol detayları). Bu iki veri yolu, HAL kütüphanesinin gerçekten değer
kattığı yerlerdir. Her ikisinin de kesme (\code{\_IT}) ve DMA (\code{\_DMA})
sürümleri vardır: örneğin \code{HAL\_SPI\_Transmit\_DMA()} ile bir ekranı CPU'yu
meşgul etmeden güncelleyebilirsin. Register seviyesini anlamak yine de hata
ayıklama için faydalıdır.}

% #####################################################################
\gpart{4}{Kesmeler, NVIC ve DMA}
% #####################################################################

% =====================================================================
\gsec{Kesmeler ve NVIC}
% =====================================================================

Şimdiye kadarki örneklerde donanımı sürekli \emph{yokladık} (polling ---
``veri geldi mi? geldi mi?''). Bu, işlemciyi meşgul eder ve verimsizdir.
\emph{Kesme} (interrupt), donanımın ``bir olay oldu!'' diye işlemciyi
\emph{uyandırmasıdır}: işlemci ne yapıyorsa duraklatır, olaya özel bir işleyici
(ISR) çalıştırır, sonra kaldığı yerden devam eder.

\gsub{NVIC: Kesme Denetleyicisi}

Cortex-M'de kesmeleri \emph{NVIC} (Nested Vectored Interrupt Controller)
yönetir. NVIC, hangi kesmelerin etkin olduğunu, önceliklerini ve iç içe
(nested) çalışmasını denetler. Her kesme kaynağının bir numarası (IRQ number)
ve bir işleyici adresi (vektör tablosunda) vardır.

\begin{center}
\begin{tabularx}{\textwidth}{@{}>{\bfseries\color{anaMavi}}l X@{}}
\toprule
\rowcolor{acikMavi}\color{anaMavi}\textbf{Kavram} & \color{anaMavi}\textbf{Açıklama}\\
\midrule
Vektör tablosu & Her kesme için işleyici fonksiyon adresini tutan tablo (Flash başında).\\
\rowcolor{acikGri} IRQ numarası & Her kesme kaynağının kimliği (ör. \code{TIM2\_IRQn}).\\
Öncelik (priority) & Küçük sayı = yüksek öncelik; yüksek olan düşüğü kesebilir.\\
\rowcolor{acikGri} ISR & Interrupt Service Routine --- kesme gelince çalışan fonksiyon.\\
\bottomrule
\end{tabularx}
\end{center}

\gsub{Timer Kesmesi Örneği}

Bir timer'ın periyodik olarak kesme üretmesini sağlayalım --- böylece
\code{delay} ile beklemek yerine, işlemci başka iş yaparken timer arka planda
LED'i çevirir.

\begin{lstlisting}
#include "stm32f1xx.h"

void timer_irq_init(void)
{
    RCC->APB1ENR |= RCC_APB1ENR_TIM2EN;

    TIM2->PSC = 7200 - 1;                  /* 10 kHz */
    TIM2->ARR = 10000 - 1;                 /* 1 Hz taşma */
    TIM2->DIER |= TIM_DIER_UIE;            /* update kesmesini etkinleştir */
    TIM2->CR1  |= TIM_CR1_CEN;

    NVIC_SetPriority(TIM2_IRQn, 2);        /* öncelik ata */
    NVIC_EnableIRQ(TIM2_IRQn);             /* NVIC'te bu kesmeyi aç */
}

/* ISR: adı vektör tablosundaki isimle TAM eşleşmeli (startup dosyasından) */
void TIM2_IRQHandler(void)
{
    if (TIM2->SR & TIM_SR_UIF) {           /* kesme kaynağını doğrula */
        TIM2->SR &= ~TIM_SR_UIF;           /* bayrağı temizle (ŞART!) */
        GPIOC->ODR ^= (1 << 13);           /* LED'i çevir */
    }
}
\end{lstlisting}

\uyari{\textbf{ISR'de kesme bayrağını temizlemeyi unutma!} Bayrağı
temizlemezsen, işleyiciden çıkar çıkmaz kesme \emph{hemen tekrar} tetiklenir ---
program sonsuza dek ISR'de kalır (kilitlenme gibi görünür). Ayrıca ISR'lerin
\emph{kısa} olması ve içinde \code{HAL\_Delay} gibi bloklayan çağrılar
bulunmaması gerekir. ISR adının startup dosyasındaki vektör adıyla birebir aynı
olması da şarttır.}

\gsub{Harici Kesme (EXTI): Buton ile}

Bir GPIO pinindeki değişimi (butona basılması) yoklamadan yakalamak için
\emph{EXTI} (External Interrupt) kullanılır. Pin durumu değiştiğinde otomatik
kesme üretilir.

\begin{lstlisting}
void exti_button_init(void)
{
    RCC->APB2ENR |= RCC_APB2ENR_IOPAEN | RCC_APB2ENR_AFIOEN;
    /* PA0'ı giriş yap (CRL) ... (giriş yapılandırması yukarıdaki gibi) */

    /* EXTI hattı 0'ı PA0'a bağla (AFIO üzerinden) */
    AFIO->EXTICR[0] &= ~AFIO_EXTICR1_EXTI0;   /* port A seç */
    EXTI->IMR  |= EXTI_IMR_MR0;                /* hat 0 kesmesini aç */
    EXTI->FTSR |= EXTI_FTSR_TR0;               /* düşen kenarda tetikle */

    NVIC_EnableIRQ(EXTI0_IRQn);
}

void EXTI0_IRQHandler(void)
{
    if (EXTI->PR & EXTI_PR_PR0) {          /* hat 0 mı tetikledi? */
        EXTI->PR = EXTI_PR_PR0;            /* pending bit'i temizle (1 yazarak) */
        GPIOC->ODR ^= (1 << 13);          /* butona basılınca LED çevir */
    }
}
\end{lstlisting}

\gsub{HAL Driver ile Harici Kesme}

HAL ile EXTI kurmak çok daha kısadır: pini \code{GPIO\_MODE\_IT\_FALLING}
moduyla ilklendirir, NVIC'i açar ve ortak \code{HAL\_GPIO\_EXTI\_Callback}
fonksiyonunu yazarsın.

\begin{lstlisting}
#include "stm32f1xx_hal.h"

void exti_hal_init(void)
{
    __HAL_RCC_GPIOA_CLK_ENABLE();

    GPIO_InitTypeDef gpio = {0};
    gpio.Pin  = GPIO_PIN_0;
    gpio.Mode = GPIO_MODE_IT_FALLING;      /* düşen kenarda kesme */
    gpio.Pull = GPIO_PULLUP;
    HAL_GPIO_Init(GPIOA, &gpio);

    HAL_NVIC_SetPriority(EXTI0_IRQn, 2, 0);
    HAL_NVIC_EnableIRQ(EXTI0_IRQn);
}

/* ISR sadece HAL'a yönlendirir */
void EXTI0_IRQHandler(void) { HAL_GPIO_EXTI_IRQHandler(GPIO_PIN_0); }

/* HAL, doğru pini bulup bu ortak callback'i çağırır */
void HAL_GPIO_EXTI_Callback(uint16_t pin)
{
    if (pin == GPIO_PIN_0) {
        HAL_GPIO_TogglePin(GPIOC, GPIO_PIN_13);
    }
}
\end{lstlisting}

% =====================================================================
\gsec{SysTick ve Zamanlama}
% =====================================================================

SysTick, her Cortex-M çekirdeğinde bulunan basit bir 24-bit geri sayım
zamanlayıcısıdır. Genelde işletim sistemleri veya \code{HAL\_Delay} için
periyodik bir ``sistem tik''i (ör. her 1 ms) üretmek üzere kullanılır ---
donanımdan bağımsız olduğu için taşınabilirdir.

\begin{lstlisting}
#include "stm32f1xx.h"

volatile uint32_t ms_ticks = 0;            /* geçen milisaniye sayacı */

void systick_init(void)
{
    /* Her 1 ms'de bir kesme: 72 MHz / 1000 = 72000 sayım */
    SysTick_Config(72000);                 /* CMSIS: yükle + kesmeyi aç */
}

/* SysTick ISR'si: her 1 ms'de bir çalışır */
void SysTick_Handler(void)
{
    ms_ticks++;
}

/* Milisaniye cinsinden bloklamalı gecikme */
void delay_ms(uint32_t ms)
{
    uint32_t start = ms_ticks;
    while ((ms_ticks - start) < ms);       /* hedef süre geçene kadar bekle */
}
\end{lstlisting}

\bilgi{SysTick tabanlı \code{delay\_ms}, meşgul döngüyle (\code{nop} sayan)
gecikmeden çok daha doğrudur: kesin milisaniye verir ve derleyici
optimizasyonundan etkilenmez. \code{ms\_ticks}'in \code{volatile} olması
kritiktir --- çünkü ISR'de değişir, ana kod ise onu okur. \code{volatile}
olmazsa derleyici \code{delay\_ms} döngüsünü sonsuza kadar sürdürebilir.}

% =====================================================================
\gsec{DMA: İşlemciyi Meşgul Etmeden Veri Taşıma}
% =====================================================================

DMA (Direct Memory Access), bellek ile çevre birimleri (veya bellekten belleğe)
arasında veriyi \emph{işlemci olmadan} taşıyan bir birimdir. Örneğin ADC'den
1000 örneği okuyup RAM'e yazmak; DMA olmadan işlemci her örneği tek tek
kopyalar (yavaş, meşgul), DMA ile ise işlemci başka iş yaparken transfer arka
planda yapılır.

\begin{center}
\begin{tabularx}{\textwidth}{@{}>{\bfseries\color{anaMavi}}l X@{}}
\toprule
\rowcolor{acikMavi}\color{anaMavi}\textbf{DMA olmadan} & \color{anaMavi}\textbf{DMA ile}\\
\midrule
İşlemci her baytı elle kopyalar & Donanım kopyalar; işlemci serbest\\
\rowcolor{acikGri} Yüksek CPU yükü & Neredeyse sıfır CPU yükü\\
Yüksek hızda veri kaybı riski & Hızlı, sürekli akış (ADC, UART, SPI)\\
\bottomrule
\end{tabularx}
\end{center}

\begin{lstlisting}
/* ADC'den belleğe DMA ile sürekli okuma (kavramsal iskelet) */
void adc_dma_init(uint16_t *buffer, uint16_t length)
{
    RCC->AHBENR |= RCC_AHBENR_DMA1EN;      /* DMA1 saatini aç */

    DMA1_Channel1->CPAR = (uint32_t)&ADC1->DR;    /* kaynak: ADC veri register'ı */
    DMA1_Channel1->CMAR = (uint32_t)buffer;       /* hedef: RAM tamponu */
    DMA1_Channel1->CNDTR = length;                /* aktarılacak eleman sayısı */
    DMA1_Channel1->CCR = DMA_CCR_MINC |           /* bellek adresini artır */
                         DMA_CCR_CIRC |           /* dairesel mod (sürekli) */
                         DMA_CCR_EN;              /* DMA kanalını başlat */
    /* Artık ADC her dönüşümü otomatik olarak buffer[]'a yazar -- CPU serbest */
}
\end{lstlisting}

\ipucu{DMA, gömülü sistemlerde yüksek verimin anahtarıdır. Sürekli ADC
örnekleme, hızlı UART/SPI akışı, ekran güncelleme gibi işlerde işlemciyi
tamamen boşa çıkarır. HAL'da \code{HAL\_ADC\_Start\_DMA()} gibi fonksiyonlar
DMA kurulumunu kolaylaştırır. Dairesel (circular) mod, aynı tamponu sürekli
döngüde doldurarak kesintisiz akış sağlar.}

% #####################################################################
\gpart{5}{Sistem Kaynakları: RTC, Watchdog ve Düşük Güç (HAL)}
% #####################################################################

STM32, çevresel iletişimin ötesinde sistem düzeyinde birçok kaynak sunar:
gerçek zaman saati (RTC), sistemi kilitlenmelerden koruyan watchdog ve pil
ömrünü uzatan düşük güç modları. Bunları HAL sürücüleriyle inceliyoruz.

% =====================================================================
\gsec{RTC: Gerçek Zaman Saati}
% =====================================================================

RTC (Real-Time Clock), tarih ve saati (yıl/ay/gün, saat/dakika/saniye) tutan
bir çevre birimidir. Ayrı bir düşük hızlı kristalle (LSE, 32.768 kHz) çalışır
ve ana güç kesilse bile bir pil (VBAT) ile saymaya devam edebilir. Veri
kaydedicilerde ve zamanlama gerektiren uygulamalarda temeldir.

\begin{lstlisting}
#include "stm32f1xx_hal.h"

RTC_HandleTypeDef hrtc;

void rtc_hal_init(void)
{
    hrtc.Instance = RTC;
    hrtc.Init.AsynchPrediv = RTC_AUTO_1_SECOND;   /* 1 saniyelik tik */
    HAL_RTC_Init(&hrtc);
}

void rtc_set_time(void)
{
    RTC_TimeTypeDef t = {0};
    t.Hours = 14; t.Minutes = 30; t.Seconds = 0;
    HAL_RTC_SetTime(&hrtc, &t, RTC_FORMAT_BIN);

    RTC_DateTypeDef d = {0};
    d.Year = 26; d.Month = 7; d.Date = 13;        /* 13 Temmuz 2026 */
    HAL_RTC_SetDate(&hrtc, &d, RTC_FORMAT_BIN);
}

void rtc_read_time(void)
{
    RTC_TimeTypeDef t;
    RTC_DateTypeDef d;
    HAL_RTC_GetTime(&hrtc, &t, RTC_FORMAT_BIN);
    HAL_RTC_GetDate(&hrtc, &d, RTC_FORMAT_BIN);    /* GetTime'dan SONRA çağır! */
    printf("%02d:%02d:%02d\n", t.Hours, t.Minutes, t.Seconds);
}
\end{lstlisting}

\uyari{RTC'de \code{HAL\_RTC\_GetTime}'dan sonra \textbf{mutlaka}
\code{HAL\_RTC\_GetDate} çağır --- saati okumak, tarih register'ını kilitler ve
tarihi okuyana kadar güncellenmez. Sırayı bozarsan tarih/saat tutarsız okunur.
Ayrıca RTC'yi yeniden başlatmamak için ilk kurulumda backup register'a bir
imza yazıp kontrol etmek yaygın bir kalıptır.}

% =====================================================================
\gsec{Watchdog: Sistemi Kilitlenmeden Koruma}
% =====================================================================

Watchdog (bekçi köpeği), yazılım kilitlenirse (sonsuz döngü, çökme) sistemi
otomatik \emph{yeniden başlatan} bir zamanlayıcıdır. Fikir: kod düzenli
aralıklarla watchdog'u ``besler'' (reset'ler); beslemeyi durdurursa (kilitlendi
demektir) watchdog süre dolunca MCU'yu resetler. Güvenlik kritik gömülü
sistemlerde şarttır.

\begin{lstlisting}
#include "stm32f1xx_hal.h"

IWDG_HandleTypeDef hiwdg;

void watchdog_hal_init(void)
{
    hiwdg.Instance       = IWDG;
    hiwdg.Init.Prescaler = IWDG_PRESCALER_64;
    hiwdg.Init.Reload    = 1250;           /* ~2 saniye zaman aşımı (40kHz LSI) */
    HAL_IWDG_Init(&hiwdg);                 /* başlar başlamaz saymaya başlar */
}

int main(void)
{
    HAL_Init();
    watchdog_hal_init();

    while (1) {
        do_work();                         /* normal iş */
        HAL_IWDG_Refresh(&hiwdg);          /* watchdog'u besle (reset'i geciktir) */
        /* Eğer do_work() kilitlenirse Refresh çağrılmaz -> MCU resetlenir */
    }
}
\end{lstlisting}

\bilgi{İki tür watchdog vardır: \textbf{IWDG} (Independent Watchdog) bağımsız
düşük hızlı saatle çalışır --- ana saat dursa bile korur, en sağlamı.
\textbf{WWDG} (Window Watchdog) ise beslemenin belirli bir ``pencere''
içinde yapılmasını ister --- hem çok geç hem çok erken beslemeyi hata sayar,
daha katı zamanlama denetimi sağlar. Kritik uygulamalarda IWDG standarttır.}

% =====================================================================
\gsec{Düşük Güç Modları}
% =====================================================================

Pil ile çalışan cihazlarda (IoT sensörü, giyilebilir) güç tüketimi kritiktir.
STM32, işi olmadığında uykuya geçerek akımı miliamperlerden mikroamperlere
düşürebilir. Üç temel mod vardır; ne kadar çok kapatırsan o kadar az güç ama
o kadar yavaş uyanış.

\begin{center}
\begin{tabularx}{\textwidth}{@{}>{\bfseries\color{anaMavi}}l X X@{}}
\toprule
\rowcolor{acikMavi}\color{anaMavi}\textbf{Mod} & \color{anaMavi}\textbf{Ne kapanır?} & \color{anaMavi}\textbf{Uyanma}\\
\midrule
Sleep & CPU durur, çevre birimleri çalışır & Herhangi bir kesme\\
\rowcolor{acikGri} Stop & CPU + saatlerin çoğu durur, RAM korunur & EXTI / RTC alarmı\\
Standby & Neredeyse her şey kapanır (en düşük güç) & Reset / WKUP pini / RTC\\
\bottomrule
\end{tabularx}
\end{center}

\begin{lstlisting}
#include "stm32f1xx_hal.h"

/* Sleep: en hafif uyku -- kesme gelene kadar CPU'yu durdur (WFI) */
void enter_sleep(void)
{
    HAL_PWR_EnterSLEEPMode(PWR_MAINREGULATOR_ON, PWR_SLEEPENTRY_WFI);
    /* Bir kesme gelince buradan devam eder */
}

/* Stop: düşük güç; EXTI (ör. buton) veya RTC alarmı ile uyanır */
void enter_stop(void)
{
    HAL_PWR_EnterSTOPMode(PWR_LOWPOWERREGULATOR_ON, PWR_STOPENTRY_WFI);
    SystemClock_Config();                  /* uyandıktan sonra saati YENİDEN kur */
}

/* Standby: en düşük güç; sadece reset/WKUP/RTC ile uyanır, RAM kaybolur */
void enter_standby(void)
{
    HAL_PWR_EnterSTANDBYMode();            /* uyanınca sistem baştan başlar */
}
\end{lstlisting}

\uyari{\textbf{Stop modundan sonra saati yeniden kur!} Stop modu PLL ve yüksek
hızlı saatleri kapattığından, uyandıktan sonra MCU varsayılan yavaş HSI ile
çalışır; \code{SystemClock\_Config()}'i tekrar çağırmazsan çevre birimleri
yanlış frekansta çalışır. Standby modunda ise cihaz uyanınca \code{main}'den
baştan başlar (RAM içeriği kaybolur) --- durumu korumak için backup
register'ları kullanılır.}

\ipucu{Tipik düşük güçlü IoT kalıbı: MCU çoğu zaman \emph{Stop} veya
\emph{Standby}'da uyur, RTC alarmı (ör. her 10 dakika) onu uyandırır, bir
sensör okur, veriyi gönderir ve tekrar uyur. Bu sayede bir pil aylarca hatta
yıllarca dayanabilir. Güç ölçümü için bir ampermetre veya ST'nin
\emph{Power Shield}'ı kullanılır.}

% =====================================================================
\gsec{CRC ve Diğer Sistem Birimleri}
% =====================================================================

STM32, veri bütünlüğünü donanımda hızlıca doğrulayan bir \emph{CRC} (Cyclic
Redundancy Check) birimi içerir. İletişim protokollerinde ve firmware
doğrulamada verinin bozulmadığını kontrol etmek için kullanılır.

\begin{lstlisting}
#include "stm32f1xx_hal.h"

CRC_HandleTypeDef hcrc;

void crc_hal_init(void)
{
    __HAL_RCC_CRC_CLK_ENABLE();
    hcrc.Instance = CRC;
    HAL_CRC_Init(&hcrc);
}

uint32_t compute_crc(uint32_t *data, uint32_t length)
{
    /* Donanım CRC'yi anında hesaplar (yazılım döngüsünden çok hızlı) */
    return HAL_CRC_Calculate(&hcrc, data, length);
}
\end{lstlisting}

\bilgi{STM32 ailesine göre başka sistem birimleri de bulunur: \textbf{DAC}
(dijital-analog çevirici --- ses/dalga üretimi), \textbf{RNG} (donanım rastgele
sayı üreteci --- kriptografi), \textbf{FSMC} (harici bellek/ekran arayüzü) ve
\textbf{option bytes} (Flash koruma, önyükleme seçimi). Hepsi CMSIS/HAL ile
erişilebilir; hangi çevre birimlerinin mevcut olduğu kartının datasheet'inde
listelenir.}

% =====================================================================
\gsec{Hata Ayıklama ve İzleme (Debugging)}
% =====================================================================

Gömülü kodda hata ayıklamak, PC'dekinden farklıdır: bir ekran veya konsol
yoktur. STM32, donanım hata ayıklama için \emph{SWD} (Serial Wire Debug)
arayüzü sunar --- sadece iki telle (SWDIO, SWCLK) kesme noktası koyma, adım
adım ilerleme ve değişken inceleme yapılır.

\begin{center}
\begin{tabularx}{\textwidth}{@{}>{\bfseries\color{anaMavi}}l X@{}}
\toprule
\rowcolor{acikMavi}\color{anaMavi}\textbf{Yöntem} & \color{anaMavi}\textbf{Kullanım}\\
\midrule
SWD + GDB & Kesme noktası, adımlama, register/bellek inceleme (en güçlü).\\
\rowcolor{acikGri} SWO / ITM & Tek telle ``printf'' izleme (trace), CPU'yu durdurmadan.\\
UART printf & Basit mesaj çıktısı (ekstra pin gerekir).\\
\rowcolor{acikGri} LED / GPIO & En ilkel ama bazen en hızlı (bir pini yakıp söndür).\\
\bottomrule
\end{tabularx}
\end{center}

\begin{lstlisting}[style=shstyle]
# OpenOCD ile GDB sunucusu başlat (ST-Link üzerinden)
$ openocd -f interface/stlink.cfg -f target/stm32f1x.cfg

# Başka terminalde GDB ile bağlan
$ arm-none-eabi-gdb firmware.elf
(gdb) target remote localhost:3333
(gdb) load                       # firmware'i karta yükle
(gdb) break main                 # main'e kesme noktası
(gdb) continue                   # çalıştır
(gdb) print my_variable          # değişken değerini gör
(gdb) step                       # tek satır ilerle
\end{lstlisting}

\gsub{SWO ile printf İzleme (ITM)}

SWO (Serial Wire Output), \code{printf} çıktısını tek bir izleme teli üzerinden
gönderir --- UART pini harcamadan ve CPU'yu neredeyse hiç yavaşlatmadan.

\begin{lstlisting}
#include "stm32f1xx.h"

/* printf'i ITM (SWO) üzerinden gönder */
int _write(int file, char *ptr, int len)
{
    for (int i = 0; i < len; i++) {
        ITM_SendChar(*ptr++);      /* CMSIS: karakteri ITM stimulus portuna yaz */
    }
    return len;
}
/* PC'de: st-link'in SWO görüntüleyicisi veya "openocd + itmdump" ile okunur */
\end{lstlisting}

\gsub{HardFault Ayıklama}

Bir program çökerse (geçersiz bellek erişimi, hizalanmamış erişim, sıfıra
bölme) çekirdek \code{HardFault\_Handler}'a atlar. Buradaki fault register'larını
inceleyerek çökme sebebini bulabilirsin.

\begin{lstlisting}
/* HardFault işleyici: çökme anındaki durumu incele */
void HardFault_Handler(void)
{
    /* SCB->CFSR: hangi fault türü (bus/usage/memory) olduğunu söyler */
    volatile uint32_t cfsr = SCB->CFSR;    /* Configurable Fault Status Reg */
    volatile uint32_t hfsr = SCB->HFSR;    /* HardFault Status Reg */
    volatile uint32_t faddr = SCB->BFAR;   /* hataya neden olan adres */
    (void)cfsr; (void)hfsr; (void)faddr;

    while (1);                             /* debugger ile buraya bakılır */
}
\end{lstlisting}

\ipucu{Bir HardFault yakaladığında, hata ayıklayıcıyla durup \code{CFSR}
register'ını (adres \code{0xE000ED28}) inceleyerek türü (bus fault, usage
fault\dots), \code{BFAR}/\code{MMFAR} ile de hatalı adresi öğrenebilirsin.
Ayrıca yığına itilmiş PC değeri (çökme anındaki komut) çökmenin \emph{tam
satırını} verir. \code{arm-none-eabi-addr2line} ile o adresi kaynak satırına
çevir.}

% =====================================================================
\gsec{FreeRTOS: Gerçek Zamanlı İşletim Sistemi}
% =====================================================================

Bir uygulama büyüdüğünde, tek bir \code{while(1)} döngüsünde her şeyi yönetmek
zorlaşır. \emph{RTOS} (Real-Time Operating System), programı bağımsız
\emph{görevlere} (task) böler ve bir \emph{zamanlayıcı} (scheduler) onları
sırayla/önceliğe göre çalıştırır --- her biri ayrı bir \code{while(1)} gibi.
FreeRTOS, STM32'de en yaygın RTOS'tur ve CubeMX ile hazır gelir.

\gsub{Görevler (Tasks)}

\begin{lstlisting}
#include "FreeRTOS.h"
#include "task.h"

/* Her görev sonsuz döngülü bir fonksiyondur, asla return etmez */
void led_task(void *params)
{
    while (1) {
        HAL_GPIO_TogglePin(GPIOC, GPIO_PIN_13);
        vTaskDelay(pdMS_TO_TICKS(500));    /* 500 ms uyu -- CPU'yu bırakır! */
    }
}

void sensor_task(void *params)
{
    while (1) {
        read_sensor();
        vTaskDelay(pdMS_TO_TICKS(100));
    }
}

int main(void)
{
    HAL_Init();
    /* ... donanım kurulumu ... */

    /* Görevleri oluştur (fonksiyon, ad, yığın boyutu, param, öncelik, handle) */
    xTaskCreate(led_task,    "LED",    128, NULL, 1, NULL);
    xTaskCreate(sensor_task, "Sensor", 256, NULL, 2, NULL);

    vTaskStartScheduler();                 /* zamanlayıcıyı başlat -- geri dönmez */
    while (1);
}
\end{lstlisting}

\notk[\code{vTaskDelay} vs meşgul bekleme]{Kritik fark: \code{vTaskDelay},
görevi uyuttuğunda CPU'yu \emph{başka görevlere} bırakır --- meşgul döngü gibi
boşa harcamaz. Zamanlayıcı, uyuyan görev beklerken diğerlerini çalıştırır. Bu
sayede birden çok iş, tek çekirdekte eşzamanlı gibi yürür.}

\gsub{Görevler Arası İletişim: Kuyruk ve Semafor}

Görevler veriyi güvenle paylaşmak için \emph{kuyruk} (queue), senkronizasyon
için \emph{semafor}/\emph{mutex} kullanır. Bunlar yarış koşullarını önler.

\begin{center}
\begin{tabularx}{\textwidth}{@{}>{\bfseries\color{anaMavi}}l X@{}}
\toprule
\rowcolor{acikMavi}\color{anaMavi}\textbf{Nesne} & \color{anaMavi}\textbf{Kullanım}\\
\midrule
Queue & Görevler/ISR arası veri aktarımı (thread-safe FIFO).\\
\rowcolor{acikGri} Semaphore & Olay bildirimi / kaynak sayımı (ör. ISR görevi uyandırır).\\
Mutex & Paylaşılan kaynağı (UART, SPI) tek görevin kullanmasını sağlar.\\
\bottomrule
\end{tabularx}
\end{center}

\begin{lstlisting}
#include "queue.h"

QueueHandle_t data_queue;

void producer_task(void *params)
{
    uint32_t value = 0;
    while (1) {
        value = read_sensor();
        xQueueSend(data_queue, &value, portMAX_DELAY);   /* kuyruğa koy */
        vTaskDelay(pdMS_TO_TICKS(100));
    }
}

void consumer_task(void *params)
{
    uint32_t received;
    while (1) {
        /* Veri gelene kadar UYU (CPU harcamaz), gelince işle */
        if (xQueueReceive(data_queue, &received, portMAX_DELAY) == pdTRUE) {
            process(received);
        }
    }
}

/* main içinde: data_queue = xQueueCreate(10, sizeof(uint32_t)); */
\end{lstlisting}

\bilgi{FreeRTOS'un ISR ile birlikte çalışan sürümleri vardır (\code{...FromISR}):
bir kesme, \code{xQueueSendFromISR} veya \code{xSemaphoreGiveFromISR} ile bir
görevi uyandırabilir. Böylece kesme işleyicisi kısa kalır (sadece bildirir),
asıl işi bir görev yapar --- Kısım 3'teki ``üst/alt yarı'' fikrinin RTOS
karşılığı. FreeRTOS, arka planda Kısım 8'de göreceğimiz \code{PendSV} kesmesiyle
görevler arası geçiş yapar.}

% #####################################################################
\gpart{6}{ARM Cortex-M Assembly Temelleri}
% #####################################################################

Şimdiye kadar STM32'yi C ile programladık. Artık bir kat aşağıya --- C
kodunun çevrildiği \emph{assembly} diline --- iniyoruz. Assembly bilmek;
startup kodunu anlamak, hata ayıklamada disassembly okumak, kritik döngüleri
optimize etmek ve donanımın gerçekte nasıl çalıştığını kavramak için gereklidir.

% =====================================================================
\gsec{ARM Mimarisi ve Programlama Modeli}
% =====================================================================

\gsub{RISC ve Load/Store Mimarisi}

ARM bir \emph{RISC} (Reduced Instruction Set Computer) mimarisidir: az sayıda,
basit ve sabit boyutlu komutları vardır. En önemli özelliği \emph{load/store}
mimarisi olmasıdır: aritmetik/mantık işlemleri \textbf{yalnızca register'lar}
üzerinde yapılır; bellek yalnızca özel \code{LDR} (load) ve \code{STR} (store)
komutlarıyla okunur/yazılır.

\begin{center}
\begin{tabularx}{\textwidth}{@{}>{\bfseries\color{anaMavi}}l X X@{}}
\toprule
\rowcolor{acikMavi}\color{anaMavi}\textbf{} & \color{anaMavi}\textbf{ARM (RISC)} & \color{anaMavi}\textbf{x86 (CISC)}\\
\midrule
Komut sayısı & Az, basit & Çok, karmaşık\\
\rowcolor{acikGri} Bellek işlemi & Sadece LDR/STR & Çoğu komut belleğe erişebilir\\
Komut boyutu & Sabit (2/4 bayt) & Değişken (1--15 bayt)\\
\rowcolor{acikGri} Güç verimi & Yüksek (gömülü için ideal) & Düşük\\
\bottomrule
\end{tabularx}
\end{center}

\gsub{Register'lar (Cortex-M Programlama Modeli)}

Cortex-M'de 16 adet 32-bit register vardır. \code{R0}--\code{R12} genel
amaçlıdır; \code{R13}--\code{R15}'in özel görevleri vardır.

\begin{center}
\begin{tabularx}{\textwidth}{@{}>{\ttfamily\bfseries\color{anaMavi}}l >{\bfseries}l X@{}}
\toprule
\rowcolor{acikMavi}\color{anaMavi}\textbf{Register} & \color{anaMavi}\textbf{Ad} & \color{anaMavi}\textbf{Görev}\\
\midrule
R0--R12 & Genel amaçlı & Veri ve ara hesaplamalar.\\
\rowcolor{acikGri} R13 & SP (Stack Pointer) & Yığının tepesini gösterir.\\
R14 & LR (Link Register) & Fonksiyon dönüş adresini tutar.\\
\rowcolor{acikGri} R15 & PC (Program Counter) & Çalışan komutun adresi.\\
--- & xPSR & Durum bayrakları (N, Z, C, V) ve durum.\\
\bottomrule
\end{tabularx}
\end{center}

\gsub{Durum Bayrakları (Condition Flags)}

Aritmetik ve karşılaştırma komutları, sonuç hakkında \emph{bayrakları}
(APSR içinde) günceller. Dallanma (branch) kararları bu bayraklara bakar.

\begin{center}
\begin{tabularx}{\textwidth}{@{}>{\ttfamily\bfseries\color{anaMavi}}l >{\bfseries}l X@{}}
\toprule
\rowcolor{acikMavi}\color{anaMavi}\textbf{Bayrak} & \color{anaMavi}\textbf{Ad} & \color{anaMavi}\textbf{Anlamı}\\
\midrule
N & Negative & Sonuç negatif (en yüksek bit 1).\\
\rowcolor{acikGri} Z & Zero & Sonuç sıfır.\\
C & Carry & Elde/ödünç oluştu (taşma, işaretsiz).\\
\rowcolor{acikGri} V & Overflow & İşaretli taşma.\\
\bottomrule
\end{tabularx}
\end{center}

\gsub{Thumb ve Thumb-2}

Klasik ARM komutları 32-bit'tir. \emph{Thumb}, aynı işleri 16-bit komutlarla
yaparak kod boyutunu küçültür. \emph{Thumb-2}, 16 ve 32-bit komutları karıştırıp
hem küçük hem güçlü olur. \textbf{Cortex-M yalnızca Thumb-2 kullanır} --- bu
yüzden STM32 assembly'sinde ARM state yoktur; her şey Thumb'dır.

\bilgi{Bu rehberde \textbf{GNU assembler} (\code{arm-none-eabi-as}) sözdizimini
kullanacağız. Dosyanın başında \code{.syntax unified} (birleşik sözdizimi) ve
\code{.thumb} (Thumb-2 modu) yönergeleri bulunur. Yorumlar \code{;} veya
\code{//} ile, etiketler \code{label:} biçiminde yazılır.}

% =====================================================================
\gsec{Temel Komutlar: Veri ve Aritmetik}
% =====================================================================

\gsub{Veri Taşıma: MOV ve LDR}

\code{MOV}, register'lar arasında veya küçük bir sabiti bir register'a taşır.
Büyük 32-bit sabitler tek komuta sığmaz; onlar için \code{MOVW}/\code{MOVT}
çifti veya \code{LDR Rd, =sabit} sözde-komutu (literal pool) kullanılır.

\begin{lstlisting}[style=asmstyle]
    MOV   r0, #5          ; r0 = 5 (küçük sabit doğrudan)
    MOV   r1, r0          ; r1 = r0 (register kopyala)
    MOVW  r2, #0x1234     ; r2'nin alt 16 biti = 0x1234
    MOVT  r2, #0x5678     ; r2'nin üst 16 biti = 0x5678 -> r2 = 0x56781234
    LDR   r3, =0x40011000 ; büyük sabiti yükle (literal pool ile) -- GPIOC adresi
\end{lstlisting}

\gsub{Bellek Erişimi: LDR ve STR}

Load/store mimarisinin kalbi: \code{LDR} bellekten register'a \emph{yükler},
\code{STR} register'dan belleğe \emph{yazar}. Bayt için \code{LDRB}/\code{STRB},
yarım kelime için \code{LDRH}/\code{STRH} kullanılır.

\begin{lstlisting}[style=asmstyle]
    LDR   r1, =myvar     ; r1 = myvar'ın ADRESİ
    LDR   r0, [r1]       ; r0 = *r1 (adresteki değeri yükle)
    ADD   r0, r0, #1     ; r0 = r0 + 1
    STR   r0, [r1]       ; *r1 = r0 (değeri belleğe geri yaz)

    LDRB  r2, [r1]       ; tek bayt yükle (alt 8 bit)
    STRB  r2, [r1, #4]   ; r1+4 adresine bayt yaz (offset ile)
\end{lstlisting}

\uyari{Aritmetik komutlar (\code{ADD}, \code{SUB}\dots) \textbf{belleğe
doğrudan erişemez}. Bir bellek konumunu artırmak üç adım gerektirir:
\code{LDR} (yükle) $\to$ \code{ADD} (değiştir) $\to$ \code{STR} (geri yaz).
Bu, x86'dan gelenler için ilk şaşırtıcı noktadır; ama load/store mimarisinin
özüdür.}

\gsub{Aritmetik ve Mantık Komutları}

\begin{center}
\begin{tabularx}{\textwidth}{@{}>{\ttfamily\bfseries\color{anaMavi}}l X@{}}
\toprule
\rowcolor{acikMavi}\color{anaMavi}\textbf{Komut} & \color{anaMavi}\textbf{İşlevi}\\
\midrule
ADD Rd, Rn, Rm & Rd = Rn + Rm (topla)\\
\rowcolor{acikGri} SUB Rd, Rn, Rm & Rd = Rn $-$ Rm (çıkar)\\
MUL Rd, Rn, Rm & Rd = Rn $\times$ Rm (çarp)\\
\rowcolor{acikGri} UDIV / SDIV & Bölme (işaretsiz / işaretli)\\
AND / ORR / EOR & Bit bazlı VE / VEYA / ÖZEL-VEYA (XOR)\\
\rowcolor{acikGri} BIC & Bit temizle (Rn AND NOT Rm)\\
LSL / LSR & Sola / sağa mantıksal kaydır\\
\bottomrule
\end{tabularx}
\end{center}

\begin{lstlisting}[style=asmstyle]
    ADD   r0, r1, r2     ; r0 = r1 + r2
    SUB   r0, r0, #10    ; r0 = r0 - 10 (sabit ile)
    MUL   r3, r1, r2     ; r3 = r1 * r2
    ADDS  r0, r1, r2     ; topla VE bayrakları güncelle (S soneki)

    ORR   r0, r0, #(1 << 13)   ; r0'ın 13. bitini kur (register set örneği)
    BIC   r0, r0, #(1 << 13)   ; r0'ın 13. bitini temizle
    LSL   r0, r1, #2           ; r0 = r1 << 2 (4 ile çarpma)
\end{lstlisting}

\notk[``S'' soneki]{\code{ADD} bayrakları değiştirmez, ama \code{ADDS}
(sondaki \code{S}) sonuç hakkında N/Z/C/V bayraklarını günceller. Bir sonraki
bölümdeki koşullu dallanmalar bu bayraklara baktığından, karşılaştırma/döngü
kodunda çoğu zaman \code{S}'li sürümler veya \code{CMP} kullanılır.}

% #####################################################################
\gpart{7}{ARM Assembly: Akış, Fonksiyonlar ve Uygulama}
% #####################################################################

% =====================================================================
\gsec{Akış Kontrolü: Dallanma ve Döngüler}
% =====================================================================

Yüksek seviyeli dillerdeki \code{if}, \code{while}, \code{for} yapıları
assembly'de \emph{karşılaştırma} + \emph{koşullu dallanma} (branch) ikilisiyle
kurulur. Önce \code{CMP} ile bayraklar ayarlanır, sonra koşullu bir \code{B}
komutu bayraklara bakarak atlar.

\gsub{Karşılaştırma ve Koşullu Dallanma}

\code{CMP r0, r1} aslında \code{r0 - r1} yapar ama sonucu \emph{saklamaz};
yalnızca bayrakları günceller. Ardından koşul ekli bir dal komutu gelir:

\begin{center}
\begin{tabularx}{\textwidth}{@{}>{\ttfamily\bfseries\color{anaMavi}}l l >{\ttfamily\bfseries\color{anaMavi}}l l@{}}
\toprule
\rowcolor{acikMavi}\color{anaMavi}\textbf{Komut} & \color{anaMavi}\textbf{Koşul} & \color{anaMavi}\textbf{Komut} & \color{anaMavi}\textbf{Koşul}\\
\midrule
BEQ & eşit ($=$) & BNE & eşit değil ($\ne$)\\
\rowcolor{acikGri} BGT & büyük ($>$, işaretli) & BLT & küçük ($<$, işaretli)\\
BGE & büyük/eşit ($\ge$) & BLE & küçük/eşit ($\le$)\\
\rowcolor{acikGri} BHI & büyük (işaretsiz) & BLO & küçük (işaretsiz)\\
B & koşulsuz (her zaman) & BX & register'a dallan (dönüş)\\
\bottomrule
\end{tabularx}
\end{center}

\gsub{if Yapısı}

\begin{lstlisting}[style=asmstyle]
    ; C karşılığı:  if (r0 > 10) { r1 = 1; } else { r1 = 0; }
    CMP   r0, #10        ; r0 ile 10'u karşılaştır (bayrakları ayarla)
    BLE   else_branch    ; r0 <= 10 ise else'e atla
    MOV   r1, #1         ; if bloğu: r1 = 1
    B     end_if         ; else'i atla
else_branch:
    MOV   r1, #0         ; else bloğu: r1 = 0
end_if:
    ; devam...
\end{lstlisting}

\gsub{Döngü: Bir Diziyi Toplama}

\begin{lstlisting}[style=asmstyle]
    ; C karşılığı:  sum = 0; for (i=0; i<n; i++) sum += array[i];
    LDR   r1, =data_array ; r1 = dizinin adresi
    MOV   r2, #0          ; r2 = sum = 0
    MOV   r3, #0          ; r3 = i = 0
    MOV   r4, #5          ; r4 = n = 5 (eleman sayısı)
loop_start:
    CMP   r3, r4          ; i < n mi?
    BGE   loop_end        ; i >= n ise döngüden çık
    LDR   r5, [r1, r3, LSL #2]  ; r5 = array[i]  (i*4 offset ile)
    ADD   r2, r2, r5      ; sum += array[i]
    ADD   r3, r3, #1      ; i++
    B     loop_start      ; döngü başına dön
loop_end:
    ; r2 artık toplamı içerir
\end{lstlisting}

\bilgi{\code{LDR r5, [r1, r3, LSL \#2]} tek bir komutta güçlü bir iş yapar:
\code{r3}'ü (indeks) 2 bit sola kaydırıp (yani 4 ile çarpıp --- \code{int}
boyutu) \code{r1}'e (taban adres) ekler ve o adresteki değeri yükler. Yani
\code{array[i]}'yi tek komutta okur. Buna \emph{ölçekli register offset}
adresleme denir --- ARM'ın verimliliğinin bir örneği.}

% =====================================================================
\gsec{Fonksiyonlar, Yığın ve Çağrı Kuralı (AAPCS)}
% =====================================================================

Fonksiyon çağırmak, \code{BL} (Branch with Link) ile yapılır: hedefe atlar ve
dönüş adresini \code{LR}'ye kaydeder. Fonksiyondan \code{BX LR} ile dönülür.
Ama fonksiyonların birlikte çalışabilmesi için herkesin uyduğu bir sözleşme
gerekir: \emph{AAPCS} (ARM Architecture Procedure Call Standard).

\gsub{Çağrı Kuralı (Calling Convention)}

\begin{center}
\begin{tabularx}{\textwidth}{@{}>{\ttfamily\bfseries\color{anaMavi}}l X@{}}
\toprule
\rowcolor{acikMavi}\color{anaMavi}\textbf{Register} & \color{anaMavi}\textbf{Rol}\\
\midrule
R0--R3 & İlk 4 argüman; ve scratch (çağıran korumaz). R0 = dönüş değeri.\\
\rowcolor{acikGri} R4--R11 & Yerel değişkenler; \textbf{callee-saved} (kullanan korumalı).\\
R12 (IP) & Scratch (ara).\\
\rowcolor{acikGri} LR (R14) & Dönüş adresi.\\
SP (R13) & Yığın işaretçisi; 8 bayt hizalı tutulmalı.\\
\bottomrule
\end{tabularx}
\end{center}

\gsub{Basit Fonksiyon (Leaf Function)}

Başka fonksiyon çağırmayan (leaf) bir fonksiyon, \code{LR}'yi saklamaya gerek
duymaz --- doğrudan \code{BX LR} ile döner.

\begin{lstlisting}[style=asmstyle]
    ; int add(int a, int b) { return a + b; }
    ; a -> r0, b -> r1, dönüş -> r0
add:
    ADD   r0, r0, r1     ; r0 = a + b (dönüş değeri r0'da)
    BX    lr             ; çağırana dön
\end{lstlisting}

\gsub{Yığın (Stack): PUSH ve POP}

Bir fonksiyon başka fonksiyon çağıracaksa, \code{BL} \code{LR}'nin üzerine
yazacağı için önce \code{LR}'yi \emph{yığına} kaydetmelidir. Ayrıca callee-saved
register'ları (R4--R11) kullanıyorsa onları da saklamalıdır. Yığın
\emph{full-descending}'dir: aşağı doğru büyür.

\begin{lstlisting}[style=asmstyle]
    ; int compute(int x) { return add(x, 10) * 2; }
compute:
    PUSH  {r4, lr}       ; r4 ve LR'yi yığına kaydet (çağrı LR'yi bozacak)
    MOV   r4, r0         ; x'i r4'te sakla (add çağrısı r0'ı bozacak)
    MOV   r1, #10
    BL    add            ; r0 = add(x, 10) -- LR değişir, r0-r3 bozulur
    LSL   r0, r0, #1     ; r0 = sonuç * 2
    POP   {r4, pc}       ; r4'ü geri yükle; LR'yi PC'ye yükleyerek DÖN
\end{lstlisting}

\notk[\code{POP \{..., pc\}} hilesi]{Dönüş adresini \code{LR}'ye geri yükleyip
sonra \code{BX LR} yapmak yerine, \code{PUSH \{..., lr\}} ile kaydedilen
dönüş adresini doğrudan \code{POP \{..., pc\}} ile \code{PC}'ye yüklemek ---
tek komutta hem register'ları geri yükler hem de fonksiyondan döner. Bu, ARM
kodunda çok yaygın ve verimli bir kalıptır.}

\gsub{Özyinelemeli Fonksiyon: Faktöriyel}

\begin{lstlisting}[style=asmstyle]
    ; int factorial(int n) { return n<=1 ? 1 : n*factorial(n-1); }
factorial:
    PUSH  {r4, lr}       ; n'i ve dönüş adresini koru
    CMP   r0, #1         ; n <= 1 mi?
    BGT   recurse        ; n > 1 ise özyinele
    MOV   r0, #1         ; taban durum: return 1
    POP   {r4, pc}
recurse:
    MOV   r4, r0         ; n'i sakla
    SUB   r0, r0, #1     ; r0 = n - 1
    BL    factorial      ; r0 = factorial(n-1)
    MUL   r0, r4, r0     ; r0 = n * factorial(n-1)
    POP   {r4, pc}       ; geri yükle ve dön
\end{lstlisting}

% =====================================================================
\gsec{Donanım, Startup ve C ile Karıştırma}
% =====================================================================

\gsub{Assembly'den GPIO Register'ına Yazma}

Kısım 2'deki LED yakmayı, şimdi doğrudan assembly ile yapalım. Sadece
register adresine bir değer yazmaktan ibarettir.

\begin{lstlisting}[style=asmstyle]
    ; GPIOC->BSRR = (1 << 13);  -- PC13 LED'ini yak (register seviyesi)
    LDR   r0, =0x40011010 ; r0 = GPIOC_BSRR adresi (BSRR ofseti 0x10)
    MOV   r1, #(1 << 13)  ; r1 = PC13 için bit maskesi
    STR   r1, [r0]        ; BSRR'ye yaz -> pin atomik olarak set edilir
\end{lstlisting}

\gsub{C ile Inline Assembly}

C kodunun içine, GCC'nin \code{asm} sözdizimiyle assembly gömebiliriz.
Kesme kapatma/açma gibi çekirdek işlemler genelde böyle yapılır.

\begin{lstlisting}
/* Cortex-M çekirdek komutlarını C'den çağırma (inline assembly) */
static inline void disable_irq(void)
{
    __asm__ volatile ("cpsid i");    /* tüm kesmeleri kapat (interrupt disable) */
}

static inline void enable_irq(void)
{
    __asm__ volatile ("cpsie i");    /* kesmeleri aç */
}

static inline void wait_for_interrupt(void)
{
    __asm__ volatile ("wfi");        /* kesme gelene kadar uyu (düşük güç) */
}
\end{lstlisting}

\gsub{Cortex-M Startup Kodu ve Vektör Tablosu}

Reset sonrası ilk çalışan kod, assembly ile yazılmış \emph{startup} kodudur.
Cortex-M'de vektör tablosunun ilk iki girdisi özeldir: başlangıç yığın
işaretçisi (SP) ve reset işleyicisinin (ilk çalışacak kod) adresi.

\begin{lstlisting}[style=asmstyle]
    .syntax unified
    .thumb

    .section .isr_vector, "a"     ; vektör tablosu (Flash başında)
    .word   _stack_top            ; [0] başlangıç SP değeri
    .word   Reset_Handler         ; [1] reset sonrası çalışacak adres
    .word   NMI_Handler           ; [2] NMI
    .word   HardFault_Handler     ; [3] hard fault
    ; ... diğer işleyiciler ...

    .section .text
    .thumb_func
Reset_Handler:
    LDR   sp, =_stack_top         ; yığın işaretçisini kur
    BL    main                    ; C main() fonksiyonuna atla
hang:
    B     hang                    ; main dönerse burada takıl (sonsuz döngü)
\end{lstlisting}

\bilgi{Gerçek startup kodu ek işler de yapar: \code{.data} bölümünü Flash'tan
RAM'e kopyalamak (ilklendirilmiş global değişkenler) ve \code{.bss} bölümünü
sıfırlamak (ilklendirilmemiş globaller). Bunlar C'nin \code{main}'den önce
beklediği ortamı hazırlar. STM32CubeIDE bu startup dosyasını (\code{startup\_stm32f103.s})
otomatik üretir; ama içinde ne olduğunu bilmek, gömülü sistemi gerçekten
anlamak demektir.}

% #####################################################################
\gpart{8}{ARM Assembly: İleri Konular}
% #####################################################################

Bu kısım, Cortex-M assembly'sinin daha güçlü özelliklerini kapsar: tek komutta
kaydırma yapan barrel shifter, zengin adresleme modları, atomik erişim, özel
register'lar, kayan nokta (FPU) ve donanım exception'larını assembly ile
işleme. Bunlar, verimli sürücüler, işletim sistemi çekirdekleri ve gerçek
zamanlı kod yazmanın anahtarıdır.

% =====================================================================
\gsec{Barrel Shifter ve Adresleme Modları}
% =====================================================================

\gsub{Barrel Shifter: Bedava Kaydırma}

ARM'ın ayırt edici özelliği \emph{barrel shifter}'dır: birçok komutun ikinci
operandı, \emph{aynı komut içinde} ekstra çevrim harcamadan kaydırılabilir.
Böylece ``4 ile çarp ve topla'' gibi işlemler tek komutta yapılır.

\begin{lstlisting}[style=asmstyle]
    ADD   r0, r1, r2, LSL #2   ; r0 = r1 + (r2 << 2)  yani r1 + r2*4
    MOV   r0, r1, LSR #1       ; r0 = r1 >> 1  (2'ye bölme)
    ADD   r0, r1, r1, LSL #3   ; r0 = r1 + r1*8 = r1*9  (çarpmasız!)
    RSB   r0, r1, r1, LSL #4   ; r0 = r1*16 - r1 = r1*15
\end{lstlisting}

\bilgi{Barrel shifter, derleyicilerin sabitle çarpmayı neden nadiren gerçek
\code{MUL} ile yaptığını açıklar: \code{x*9} yerine
\code{ADD r0, r1, r1, LSL \#3} tek çevrimde biter. Kaydırma türleri: \code{LSL}
(mantıksal sol), \code{LSR} (mantıksal sağ), \code{ASR} (aritmetik sağ ---
işaret korur), \code{ROR} (döndür).}

\gsub{Adresleme Modları}

\code{LDR}/\code{STR} komutlarının adres hesaplama biçimleri (addressing modes)
zengindir. Özellikle \emph{ön/son indeksleme}, bir işaretçiyi ilerletirken
belleğe erişmeyi tek komutta yapar --- dizi/tampon döngülerinde çok verimli.

\begin{center}
\begin{tabularx}{\textwidth}{@{}>{\ttfamily\bfseries\color{anaMavi}}l X@{}}
\toprule
\rowcolor{acikMavi}\color{anaMavi}\textbf{Mod} & \color{anaMavi}\textbf{Anlamı}\\
\midrule
{[r1]} & Adres = r1 (doğrudan).\\
\rowcolor{acikGri} {[r1, \#4]} & Adres = r1 + 4 (offset; r1 değişmez).\\
{[r1, r2]} & Adres = r1 + r2 (register offset).\\
\rowcolor{acikGri} {[r1, r2, LSL \#2]} & Adres = r1 + r2*4 (ölçekli --- array[i]).\\
{[r1, \#4]!} & Ön-indeks: r1 += 4, sonra eriş (r1 güncellenir).\\
\rowcolor{acikGri} {[r1], \#4} & Son-indeks: eriş, sonra r1 += 4.\\
\bottomrule
\end{tabularx}
\end{center}

\begin{lstlisting}[style=asmstyle]
    ; Bir diziyi son-indeksleme ile gez (işaretçi otomatik ilerler)
    LDR   r1, =buffer        ; r1 = tampon başlangıcı
copy_loop:
    LDR   r0, [r1], #4       ; r0 = *r1; sonra r1 += 4 (bir sonraki word)
    ; ... r0'ı işle ...
    ; döngü koşulu ile copy_loop'a dön
\end{lstlisting}

\gsub{LDM / STM: Çoklu Yükleme/Saklama}

Tek komutla birden çok register'ı arka arkaya belleğe yazabilir veya
okuyabiliriz. Bu, blok kopyalama ve yığın işlemlerinin (\code{PUSH}/\code{POP}
aslında \code{STMDB}/\code{LDMIA}'dır) temelidir.

\begin{lstlisting}[style=asmstyle]
    ; 8 word'ü tek komutta kopyala (blok transfer)
    LDMIA r0!, {r4-r7}       ; r0'dan 4 word yükle -> r4,r5,r6,r7; r0 += 16
    STMIA r1!, {r4-r7}       ; r1'e 4 word yaz; r1 += 16
    ; IA = Increment After (adres arttıkça); PUSH = STMDB sp!, POP = LDMIA sp!
\end{lstlisting}

% =====================================================================
\gsec{Koşullu Yürütme ve Bit Komutları}
% =====================================================================

\gsub{IT Blokları (If-Then)}

Thumb-2'de, dallanma yapmadan bir-dört komutu koşullu çalıştırmak için
\emph{IT} (If-Then) bloğu kullanılır. Kısa \code{if} yapılarında dal
maliyetinden (pipeline flush) kaçınır.

\begin{lstlisting}[style=asmstyle]
    ; if (r0 == 0) r1 = 1; else r1 = 2;
    CMP   r0, #0
    ITE   EQ                 ; If-Then-Else: sonraki EQ, ondan sonraki NE
    MOVEQ r1, #1             ; r0 == 0 ise çalışır
    MOVNE r1, #2             ; r0 != 0 ise çalışır
\end{lstlisting}

\gsub{Bit Alanı ve Sıralama Komutları}

Cortex-M, bit manipülasyonu için özel komutlar sunar --- gömülü register
işlemede çok kullanışlıdır.

\begin{center}
\begin{tabularx}{\textwidth}{@{}>{\ttfamily\bfseries\color{anaMavi}}l X@{}}
\toprule
\rowcolor{acikMavi}\color{anaMavi}\textbf{Komut} & \color{anaMavi}\textbf{İşlevi}\\
\midrule
UBFX / SBFX & Bir bit alanını çıkar (unsigned/signed bit field extract).\\
\rowcolor{acikGri} BFI / BFC & Bit alanı ekle / temizle (insert/clear).\\
CLZ & Baştaki sıfır bitleri say (count leading zeros).\\
\rowcolor{acikGri} RBIT & Bit sırasını tersine çevir.\\
REV / REV16 & Bayt sırasını tersle (endian dönüşümü).\\
\bottomrule
\end{tabularx}
\end{center}

\begin{lstlisting}[style=asmstyle]
    UBFX  r0, r1, #4, #3     ; r1'in 4. bitinden başlayan 3 biti r0'a çıkar
    BFI   r1, r0, #8, #4     ; r0'ın alt 4 bitini r1'in 8. bitine yerleştir
    CLZ   r0, r1             ; r1'de baştaki sıfır sayısı (öncelik kodlama)
    REV   r0, r1             ; bayt sırasını tersle (big<->little endian)
\end{lstlisting}

% =====================================================================
\gsec{Özel Register'lar, Atomik Erişim ve Bariyerler}
% =====================================================================

\gsub{Özel Register Erişimi: MRS ve MSR}

Genel register'ların dışındaki \emph{özel} register'lara (durum, kesme maskesi,
yığın işaretçileri) yalnızca \code{MRS} (oku) ve \code{MSR} (yaz) komutlarıyla
erişilir. RTOS ve çekirdek kodu bunları yoğun kullanır.

\begin{center}
\begin{tabularx}{\textwidth}{@{}>{\ttfamily\bfseries\color{anaMavi}}l X@{}}
\toprule
\rowcolor{acikMavi}\color{anaMavi}\textbf{Register} & \color{anaMavi}\textbf{İşlevi}\\
\midrule
PRIMASK & Tüm kesmeleri kapat/aç (cpsid/cpsie i bunu değiştirir).\\
\rowcolor{acikGri} BASEPRI & Belirli öncelik altındaki kesmeleri maskele.\\
CONTROL & Ayrıcalık seviyesi ve MSP/PSP seçimi.\\
\rowcolor{acikGri} MSP / PSP & Main / Process yığın işaretçileri (RTOS ayrımı).\\
\bottomrule
\end{tabularx}
\end{center}

\begin{lstlisting}[style=asmstyle]
    MRS   r0, PRIMASK        ; mevcut kesme maskesini oku (kaydet)
    CPSID i                  ; tüm kesmeleri kapat (kritik bölge)
    ; ... kritik kod ...
    MSR   PRIMASK, r0        ; önceki maske durumunu geri yükle

    MRS   r0, PSP            ; process yığın işaretçisini oku (RTOS bağlam değiştirme)
\end{lstlisting}

\gsub{Atomik Erişim: LDREX ve STREX}

Çok görevli veya kesme paylaşımlı kodda ``oku-değiştir-yaz'' atomik olmalıdır.
\code{LDREX}/\code{STREX} çifti bunu sağlar: \code{LDREX} özel bir işaretle
okur, \code{STREX} yalnızca arada kimse yazmadıysa başarılı olur. RTOS
kilitlerinin (mutex, semafor) temelidir.

\begin{lstlisting}[style=asmstyle]
    ; Atomik artırma: *r0 += 1 (kesintiye dayanıklı)
atomic_inc:
    LDREX r1, [r0]           ; değeri özel (exclusive) olarak oku
    ADD   r1, r1, #1         ; artır
    STREX r2, r1, [r0]       ; geri yaz; r2=0 başarı, r2=1 başarısız
    CMP   r2, #0             ; başarısızsa (araya girildiyse)...
    BNE   atomic_inc         ; ...baştan dene
    BX    lr
\end{lstlisting}

\gsub{Bellek Bariyerleri: DMB, DSB, ISB}

Cortex-M işlemcileri bellek erişimlerini yeniden sıralayabilir. Belirli
sıralamanın garanti edilmesi gereken yerlerde (çevre birimi yapılandırması,
bağlam değiştirme) \emph{bellek bariyerleri} kullanılır.

\begin{center}
\begin{tabularx}{\textwidth}{@{}>{\ttfamily\bfseries\color{anaMavi}}l X@{}}
\toprule
\rowcolor{acikMavi}\color{anaMavi}\textbf{Bariyer} & \color{anaMavi}\textbf{Garantisi}\\
\midrule
DMB & Data Memory Barrier: önceki bellek erişimleri sonrakilerden önce görünür.\\
\rowcolor{acikGri} DSB & Data Sync Barrier: önceki tüm erişimler tamamlanmadan devam etmez.\\
ISB & Instruction Sync Barrier: pipeline'ı temizler (yeni ayarları uygula).\\
\bottomrule
\end{tabularx}
\end{center}

\notk[Ne zaman gerekir?]{Günlük C kodunda nadiren; ama bir çevre biriminin
saatini açtıktan hemen sonra register'ına yazarken, veya vektör tablosunu
değiştirdikten sonra \code{DSB}/\code{ISB} gerekebilir. CMSIS bunları
\code{\_\_DMB()}, \code{\_\_DSB()}, \code{\_\_ISB()} intrinsic'leriyle sunar ---
assembly yazmadan C'den çağırabilirsin.}

% =====================================================================
\gsec{FPU ve Assembly'de Exception İşleme}
% =====================================================================

\gsub{Kayan Nokta (FPU) Komutları}

Cortex-M4F ve M7 çekirdekleri bir donanım \emph{FPU} (Floating-Point Unit)
içerir; kayan nokta işlemlerini yazılım öykünmesinden çok daha hızlı yapar.
FPU komutları \code{V} ile başlar ve ayrı \code{S0}--\code{S31} register'larını
kullanır.

\begin{lstlisting}[style=asmstyle]
    VLDR  s0, [r0]          ; bellekten float yükle -> s0
    VLDR  s1, [r1]
    VADD.F32 s2, s0, s1     ; s2 = s0 + s1 (tek hassasiyetli float toplama)
    VMUL.F32 s2, s2, s0     ; s2 = s2 * s0
    VSTR  s2, [r2]          ; sonucu belleğe yaz
    VCVT.S32.F32 r3, s2     ; float -> int dönüşümü
\end{lstlisting}

\uyari{FPU'yu kullanmadan önce \textbf{etkinleştirmelisin}: \code{CPACR}
register'ında CP10/CP11 erişimini açman gerekir (genelde startup kodu yapar).
Ayrıca FPU register'ları da bağlam (context) parçasıdır; bir RTOS veya kesme
FPU kullanıyorsa, bağlam değiştirmede \code{S0-S31}'in de kaydedilmesi gerekir
--- Cortex-M bunun için ``lazy stacking'' mekanizması sunar.}

\gsub{Assembly ile Exception İşleyici}

Cortex-M'de bir kesme/exception oluştuğunda, donanım \emph{otomatik olarak}
sekiz register'ı (R0-R3, R12, LR, PC, xPSR) yığına iter ve LR'ye özel bir
\code{EXC\_RETURN} değeri koyar. Bu sayede ISR'ler \emph{normal C fonksiyonu}
gibi yazılabilir. Assembly ile bir işleyici yazmak, bağlam değiştirme
(context switch) gibi işler için gerekir.

\begin{lstlisting}[style=asmstyle]
    .thumb_func
PendSV_Handler:              ; RTOS bağlam değiştirme genelde burada olur
    MRS   r0, PSP            ; mevcut görevin yığın işaretçisini al
    STMDB r0!, {r4-r11}      ; kalan register'ları (donanım R0-R3'ü zaten itti) kaydet
    ; ... görev yığın işaretçisini kaydet, sonraki görevi yükle ...
    LDMIA r0!, {r4-r11}      ; yeni görevin register'larını geri yükle
    MSR   PSP, r0
    BX    lr                 ; EXC_RETURN ile normal göreve dön (donanım gerisini yükler)
\end{lstlisting}

\bilgi{Donanımın otomatik ``stacking''i (R0-R3, R12, LR, PC, xPSR'yi itmesi),
ISR'leri C'de yazmayı mümkün kılar --- çünkü çağıran-korumalı (caller-saved)
register'lar zaten kaydedilmiştir. \code{EXC\_RETURN} (LR'deki $0xFFFF\_FFFx$
değeri), işlemciye hangi yığına (MSP/PSP) döneceğini ve FPU bağlamı olup
olmadığını söyler. Bu mekanizma, FreeRTOS gibi çekirdeklerin \code{PendSV}
ile görevler arası geçiş yapmasının temelidir.}

% =====================================================================
\gsec{Kapanış: Gömülü Geliştirme Yol Haritası}
% =====================================================================

Bu rehberde STM32'yi register seviyesinden başlayıp assembly'ye kadar iki yönde
inceledik. İki bakış birbirini tamamlar: C ile hızlı geliştirir, assembly ile
altta ne olduğunu anlarsın.

\begin{center}
\begin{tabularx}{\textwidth}{@{}>{\bfseries\color{anaMavi}}l X@{}}
\toprule
\rowcolor{acikMavi}\color{anaMavi}\textbf{İhtiyaç} & \color{anaMavi}\textbf{Kullanılan mekanizma}\\
\midrule
Pin sürme/okuma & GPIO register'ları (ODR, IDR, BSRR)\\
\rowcolor{acikGri} Çevre birimini açma & RCC saat etkinleştirme\\
Kesin zamanlama / PWM & Timer (PSC, ARR, CCR)\\
\rowcolor{acikGri} PC ile iletişim & UART\\
Analog okuma & ADC\\
\rowcolor{acikGri} Sensör veri yolu & SPI / I2C\\
Olay tabanlı tepki & Kesme + NVIC (yoklama yerine)\\
\rowcolor{acikGri} CPU'suz veri taşıma & DMA\\
Altyapıyı anlama / optimizasyon & ARM Cortex-M assembly\\
\bottomrule
\end{tabularx}
\end{center}

\ipucu{Öğrenmenin en iyi yolu \textbf{gerçek donanımla pratik yapmaktır}.
Ucuz bir ``Blue Pill'' veya Nucleo kartı al, önce LED yak, sonra buton oku,
UART ile PC'ye mesaj gönder, bir kesme kur. Her adımda \code{arm-none-eabi-objdump
-d} ile C kodunun ürettiği assembly'yi incele --- bu, iki kısmı birbirine bağlar.
Başvuru için: STM32 \emph{Reference Manual} (register detayları),
\emph{Cortex-M Generic User Guide} ve \emph{ARMv7-M Architecture Reference
Manual} (assembly). Bol dene, boz, hata ayıkla --- gömülü sistemi öğrenmenin
tek yolu budur. İyi kodlamalar!}

\vfill
\begin{center}
{\color{ortaGri}\rule{0.5\linewidth}{0.6pt}}\\[6pt]
{\color{ortaGri}\small Rehberin sonu --- Register'lardan komutlara, iyi gömülü geliştirmeler!}
\end{center}

\end{document}
